\documentclass[12pt]{article}
\pdfoutput=1


% Use Chancery Font
\DeclareFontFamily{OT1}{pzc}{}
\DeclareFontShape{OT1}{pzc}{m}{it}{<-> s * [1.10] pzcmi7t}{}
\DeclareMathAlphabet{\mathpzc}{OT1}{pzc}{m}{it}
\newcommand{\josh}[1]{\textcolor{blue}{[Joshua: #1]}}
\newcommand{\cw}[1]{\textcolor{gray}{[Charles: #1]}}
%\usepackage{chngcntr}
%\counterwithout{equation}{section}



\newcommand{\bb}[1]{\mathbb{#1}}
\newcommand{\wt}[1]{\widetilde{#1}}
\newcommand{\ol}[1]{\overline{#1}}
\newcommand{\mc}[1]{\mathcal{#1}}
\newcommand{\hmch}{\hat{\mc{H}}}
\newcommand{\citeme}[1][]{{\color{red}[*#1]}}

\usepackage{stmaryrd}


\usepackage{draft}
\usepackage[weather]{ifsym}

\usepackage{hyperref}
\usepackage{graphicx,color,subfig}
\usepackage{cite}
\usepackage{mciteplus}
\usepackage{skak}
\usepackage{empheq}
\usepackage{tikz}
\usepackage{bbm}

% Use Chancery Font
\DeclareFontFamily{OT1}{pzc}{}
\DeclareFontShape{OT1}{pzc}{m}{it}{<-> s * [1.10] pzcmi7t}{}
\DeclareMathAlphabet{\mathpzc}{OT1}{pzc}{m}{it}

\usetikzlibrary{calc}
\usetikzlibrary{snakes}
\usetikzlibrary{arrows.meta}
\usetikzlibrary{decorations.pathmorphing}
\usetikzlibrary{decorations.markings}
\usetikzlibrary{bending}
\tikzset{snake it/.style={decorate, decoration=snake}}
\usetikzlibrary{shapes.misc}
\tikzset{cross/.style={cross out, draw=black, minimum size=2*(#1-\pgflinewidth), inner sep=0pt, outer sep=0pt},
%default radius will be 1pt.
cross/.default={1pt}}

\usepackage[T1]{fontenc}
\usepackage{esint}
\usepackage{lmodern}

\newcommand{\fixme}[1]{{\bf {\color{red}[#1]}}}
\newcommand{\vv}[1]{\left\langle #1 \right\rangle}
\newcommand{\un}[1]{\underline{#1}}
\newcommand{\BL}[1]{{ {\color{blue}[#1]}}}

\newcommand{\bgcom}[1]{\fixme{BG: #1}}

\def\be#1\ee{\begin{align}#1\end{align}}
\newcommand\nn{\nonumber}

\newcommand{\calA}{\mathcal{A}}
\newcommand{\bz}{\bar{z}}

\definecolor{dark green}{rgb}{0.7,1,0.64}

\usepackage{listings}
\usepackage{xcolor}

\definecolor{codegreen}{rgb}{0,0.6,0}
\definecolor{codegray}{rgb}{0.5,0.5,0.5}
\definecolor{codepurple}{rgb}{0.58,0,0.82}
\definecolor{backcolour}{rgb}{0.95,0.95,0.92}

\lstdefinestyle{myStyle}{
    belowcaptionskip=1\baselineskip,
    breaklines=true,
    frame=none,
    numbers=none,
    basicstyle=\footnotesize\ttfamily,
    keywordstyle=\bfseries\color{green!40!black},
    commentstyle=\itshape\color{purple!40!black},
    identifierstyle=\color{blue},
    backgroundcolor=\color{gray!10!white},
    tabsize=2,
}



\lstset{style=myStyle}
\usepackage{array}

\begin{document}

\unitlength = .8mm


\section{Heterotic worldsheet theory of standard embedding}

We begin with the heterotic worldsheet theory based on a $c=9$ $(2,2)$ SCFT combined with 26 chiral fermions $\lambda^{A'}$, 4 free bosons $X^\mu$ and 4 anti-chiral free fermions $\wt\psi^\mu$, along with $b,c,\wt b,\wt c, \wt \B,\wt \C$ ghost system. Let $(q, \wt q)$ be the (left,right) R-charge of the $(2,2)$ SCFT (integer valued in the NS sector and half-integer valued in the R sector), and $(F, \wt F)$ be the fermions numbers of $\lambda^{A'}$ and $\wt\psi^\mu,\wt\B,\wt\C$ respectively (similarly half-integer valued in the R sector). The GSO projection for the standard embedding of $SO(32)$ heterotic string is
\ie
F+q, ~\wt F + \wt q\in 2 \mathbb{Z}.
\fe
For $E_8\times E_8$ heterotic string, we impose the GSO projection to the $(2,2)$ SCFT combined with 10 chiral fermions, and separately the GSO projection on the remaining 16 chiral fermions to give the $(E_8)_1$ chiral CFT.

In particular, the $U(1)$ gauge symmetry of the 4D effective theory is associated with the $U(1)_R$ current of the holomorphic ${\cal N}=2$ superconformal algebra on the worldsheet. The charged massless fields organized into 4D ${\cal N}=1$ chiral multiplets, consisting of complex scalars and chiral fermions represented by NS and R sector vertex operators of the form
\fixme{The operators below are in the 10 and 1 under decomposition $27 \to 16 \oplus 10 \oplus 1$ of $E_6$ under $SO(10)$, the 16 is in the R sector. (2.3) is then $27^3$ superpotential}
\ie\label{nsrmasslfrbsf}
& e^{-{\wt\phi}}    {\cal O}^{\rm NS} \lambda^{A'} e^{ik\cdot X},~~~~e^{-{\wt\phi\over 2}} \wt S_{\A} {\cal O}^{\rm R} \lambda^{A'} e^{ik\cdot X},
\\
&  {\rm and}~~~ e^{-{\wt\phi}} {\cal O}'^{\rm NS} e^{ik\cdot X},~~~~ e^{-{\wt\phi\over 2}} \wt S_{\A} {\cal O}'^{\rm R} e^{ik\cdot X},
\fe
where $\wt S_\A$ is a chiral spin field of the $\mathbb{R}^{1,3}$ sector, ${\cal O}^{\rm NS}$ is a chiral primary of the $(2,2)$ SCFT with weight $({1\over 2}, {1\over 2})$ and R-charge $(q,\wt q) = (\pm 1, 1)$, and ${\cal O}^{\rm R}$ is an (NS,R) sector operator of the $(2,2)$ SCFT with weight $({1\over 2},{3\over 8})$ and R-charge $(q,\wt q) = (\pm 1, -{1\over 2})$, related to ${\cal O}^{\rm NS}$ by half unit anti-holomorphic spectral flow, i.e.
\ie\label{omrspecfl}
{\cal O}^{\rm R} = e^{-{i\over 2} \sqrt{3} \wt \varphi} {\cal O}^{\rm NS},
\fe
where $\wt\varphi$ is an anti-chiral boson defined as
\ie\label{thetaggsp}
\wt J \simeq i \sqrt{3} \bar\partial\wt \varphi,~~~~ \wt {\bf X}^\pm \simeq e^{\pm i \sqrt{3} \wt \varphi},~~~~ \wt \Theta_\pm \simeq e^{ \pm {i\over 2} \sqrt{3} \wt \varphi },
\fe
and the RHS of (\ref{omrspecfl}) is understood as normal ordered with respect to $\wt\varphi$ (which also appears in ${\cal O}^{\rm NS}$). ${\cal O}'^{\rm NS}$ is a chiral primary of the $(2,2)$ SCFT of weight $(1, {1\over 2})$ and R-charge $(q,\wt q) = (\mp 2, 1)$, related to ${\cal O}^{\rm NS}$ by $\mp 1$ unit of holomorphic spectral flow, and similarly ${\cal O}'^{\rm R}$ is an (NS,R) sector operator of weight $(1, {3\over 8})$ and R-charge $(q,\wt q) = (\mp 2, -{1\over 2})$.

There are also uncharged massless fields involving $G_{-{1\over 2}}^\pm$ descendants of $(2,2)$ chiral primaries, of the form
\ie\label{nsrmasslfrbsfneutral}
e^{-{\wt\phi}}  G_{-{1\over 2}}^\pm  {\cal O}^{{\rm NS},\mp} e^{ik\cdot X},~~~~e^{-{\wt\phi\over 2}} \wt S_{\A} G_{-{1\over 2}}^\pm  {\cal O}^{{\rm R},\mp} e^{ik\cdot X},
\fe
where the superscript $\mp$ on ${\cal O}^{\rm NS}$ and ${\cal O}^{\rm R}$ indicate their holomorphic R-charge $q=\mp 1$.

\section{4D ${\cal N}=1$ Superpotential}

The cubic superpotential terms, or equivalently the Yukawa couplings, are computed by the 3-point functions of an NS vertex operator of the form
\ie{}
e^{-{\wt\phi}}  {\cal O}^{\rm NS} \lambda^{A'}, ~~~~ e^{-{\wt\phi}}  G_{-{1\over 2}}^\pm  {\cal O}^{{\rm NS},\mp}, ~~~~ e^{-{\wt\phi}} {\cal O}'^{\rm NS} ,
\fe
and a pair of R vertex operators of the form
\ie
e^{-{\wt\phi\over 2}} \wt S_{\A} {\cal O}^{\rm R} \lambda^{A'},~~~~ e^{-{\wt\phi\over 2}} \wt S_{\A} G_{-{1\over 2}}^\pm  {\cal O}^{{\rm R},\mp}  ,~~~~  e^{-{\wt\phi\over 2}} \wt S_{\A} {\cal O}'^{\rm R}.
\fe
An example of a potentially non-vanishing 3-point correlator as constrained by global symmetries is
\ie{}
& \left\langle e^{-{\wt\phi}} {\cal O}{}'^{{\rm NS},(\pm 2, 1)} e^{-{\wt\phi\over 2}} \wt S_{\A} {\cal O}{}^{{\rm R},(\mp 1, -{1\over 2})} \lambda^{A'} e^{-{\wt\phi\over 2}} \wt S_{\B} {\cal O}{}^{{\rm R},(\mp 1, -{1\over 2})} \lambda^{B'} \right\rangle
\\
& = \epsilon_{\A\B}\delta^{A'B'} \left\langle {\cal O}{}'^{{\rm NS},(\pm 2, 1)} {\cal O}{}^{{\rm R},(\mp 1, -{1\over 2})} {\cal O}{}^{{\rm R},(\mp 1, -{1\over 2})} \right\rangle^{(2,2)~{\rm SCFT}}.
\fe
Another example of such a 3-point function is
\ie
\left\langle {\cal O}^{{\rm NS},(\pm 1, 1)} {\cal O}{}^{{\rm R},(\pm 1, -{1\over 2})} {\cal O}'{}^{{\rm R},(\mp 2, -{1\over 2})} \right\rangle^{(2,2)~{\rm SCFT}},
\fe
which is related by spectral flow or spacetime supersymmetry, and corresponds to the same cubic superpotential term.

One may also consider the 3-point function
\ie{}
& \left\langle e^{-{\wt\phi}} G_{-{1\over 2}}^\pm {\cal O}{}^{{\rm NS},(\mp 1, 1)} e^{-{\wt\phi\over 2}} \wt S_{\A} {\cal O}{}^{{\rm R},(\pm 1, -{1\over 2})} \lambda^{A'} e^{-{\wt\phi\over 2}} \wt S_{\B} {\cal O}{}^{{\rm R},(\mp 1, -{1\over 2})} \lambda^{B'} \right\rangle
\\
& = \epsilon_{\A\B}\delta^{A'B'} \left\langle G_{-{1\over 2}}^\pm {\cal O}{}^{{\rm NS},(\pm 1, 1)} {\cal O}{}^{{\rm R},(\pm 1, -{1\over 2})} {\cal O}{}^{{\rm R},(\mp 1, -{1\over 2})} \right\rangle^{(2,2)~{\rm SCFT}}
\fe
which satisfies R-charge conservation. However, the correlator in question actually vanishes due to the Ward identities associated with the holomorphic ${\cal N}=2$ SCA. This can be seen by representing $G_{-{1\over 2}}^\pm$ acting on the chiral primary at $z=\infty$ as $\oint {dz\over 2\pi i} f(z) G^\pm(z)$, where $f(z) = z + a_0  + a_1 z^{-1} + \cdots$ as $z\to \infty$. In particular we can choose $f(z) = z$, so that by shrinking the contour we obtain $G^\pm_{1\over 2}$ acting on the operator at $z=0$ and $G^\pm_{-{1\over 2}}$ acting on the operator at $z=1$, and the result vanishes as the operator at $z=0$ is a chiral primary of R-charge $q=\pm 1$. This conclusion is also consistent with the selection rule in section 2.2(c) of \cite{Lin:2016gcl}.

So we learn that the cubic superpotential is governed by the chiral ring of the $(2,2)$ SCFT. We now focus on the flat directions with respect to the cubic superpotential, and ask if there are higher order obstructions to deforming the vacuum.

\section{Deformations away from the standard embedding locus}

We view a deformation of the 4D vacuum as a worldsheet perturbation by an integrated, BRST-invariant vertex operator ${\cal V}$,
\ie\label{wsdeformation}
S \;\to\; S + \lambda \int d^2z~{\cal V}(z,\bz),
\fe
with ${\cal V}$ of total weight $(1,1)$ and compatible with the GSO projection. At first order in $\lambda$, such perturbations are in correspondence with the massless scalar vertex operators reviewed in (\ref{nsrmasslfrbsf})--(\ref{nsrmasslfrbsfneutral}). Exact marginality is the statement that the perturbation (\ref{wsdeformation}) can be continued to all orders in $\lambda$ without generating a beta function, equivalently that the corresponding spacetime scalar remains a flat direction after including all tree-level (and eventually loop) corrections.

The standard embedding locus is characterized by the factorized worldsheet theory built from a $c=9$ internal $(2,2)$ SCFT and a free left-moving gauge sector, together with the standard GSO projection. Moving away from this locus in the heterotic moduli problem means turning on deformations that preserve the anti-holomorphic ${\cal N}=2$ algebra responsible for 4D ${\cal N}=1$ supersymmetry, but break the holomorphic ${\cal N}=2$ structure of the internal SCFT. In the conventional $(\text{left},\text{right})$ notation this is a flow from $(2,2)$ to $(0,2)$ worldsheet supersymmetry (sometimes denoted $(2,0)$ in the opposite convention).

\subsection{Complete classification of supersymmetric marginal perturbations}

In the heterotic theory, preserving 4D ${\cal N}=1$ supersymmetry requires preserving the anti-holomorphic ${\cal N}=2$ algebra on the worldsheet. A convenient way to parametrize such supersymmetric deformations is to write them as right-moving F-term perturbations
\ie\label{rFterm}
S \;\to\; S + \sum_i \lambda^i \int d^2z~\Big(\wt G^-_{-{1\over 2}}{\cal W}_i(z,\bz) + \wt G^+_{-{1\over 2}}\overline{\cal W}_i(z,\bz)\Big),
\fe
where ${\cal W}_i$ are anti-holomorphic chiral primaries of the undeformed theory with
\ie
(\wt h,\wt q) = \Big({1\over 2},\,1\Big), \qquad \wt G^+_{-{1\over 2}}{\cal W}_i = 0.
\fe
Equivalently, in the unintegrated $(-1)$ picture the corresponding massless scalar vertex operators are of the form $e^{-{\wt\phi}}{\cal W}_i$ (at zero momentum), so classifying supersymmetric marginal perturbations reduces to classifying the operators ${\cal W}_i$.

At a standard embedding point the worldsheet matter CFT factorizes into the internal $(2,2)$ SCFT and a free left-moving gauge-fermion system generated by $\lambda^{A'}$. Restricting to 4D Lorentz-invariant scalar deformations, the relevant $(-1)$ picture operators ${\cal W}_i$ involve only the internal SCFT and the $\lambda^{A'}$ system (i.e.\ we exclude insertions such as $\partial X^\mu$ or $\wt\psi^\mu$). Since the theory is a tensor product, it suffices to classify factorized terms
\ie
{\cal W} \;=\; \Lambda(z)\,\mc{O}(z,\bz),
\fe
where $\Lambda$ is an operator in the free $\lambda^{A'}$ system and $\mc{O}$ is an operator in the internal $(2,2)$ SCFT.

The conditions for ${\cal W}$ to generate a supersymmetric marginal perturbation are:
\begin{enumerate}
\item {\bf Right chirality:} $\mc{O}$ is a right-moving chiral primary with $(\wt h,\wt q)=({1\over 2},1)$ (so that $\wt G^-_{-{1\over 2}}\mc{O}$ has $\wt h=1$).
\item {\bf Left weight:} $h(\Lambda)+h(\mc{O})=1$.
\item {\bf GSO:} $F(\Lambda)+q(\mc{O})\in 2\bb{Z}$, where $F(\Lambda)$ is the $\lambda^{A'}$ fermion number and $q(\mc{O})$ is the holomorphic $U(1)_R$ charge of the internal SCFT.
\end{enumerate}

\subsubsection*{Gauge-fermion sector: the only possibilities}

In the NS sector of a free fermion system, the modes are $\lambda^{A'}_{-r}$ with $r\in \bb{Z}+{1\over 2}$, contributing conformal weight $r$. Hence any state created by $F$ fermions has weight at least $F/2$. Since $h(\Lambda)\le 1$ in (\ref{rFterm}), we must have $F(\Lambda)\le 2$. Up to normalization (and dropping total derivatives, which have weight $\ge {3\over 2}$), the only independent choices are:
\begin{enumerate}
\item $\Lambda=\mathbf{1}$ with $(h,F)=(0,0)$,
\item $\Lambda=\lambda^{A'}$ with $(h,F)=({1\over 2},1)$,
\item $\Lambda=J^{A'B'}\equiv :\lambda^{A'}\lambda^{B'}:$ with $(h,F)=(1,2)$ (a gauge current).
\end{enumerate}
This proves that there are no further possibilities coming from the $\lambda^{A'}$ system at the level of massless scalar deformations.

\subsubsection*{Resulting classification}

We now combine the three gauge-fermion possibilities with the internal ${\cal N}=2$ unitarity bound $h\ge {|q|\over 2}$.
\begin{enumerate}
\item {\bf One gauge fermion ($\Lambda=\lambda^{A'}$): charged matter directions.}\\
Here $h(\mc{O})={1\over 2}$ and the GSO condition forces $q(\mc{O})$ to be odd. The unitarity bound then implies $q(\mc{O})=\pm 1$ and $h={|q|\over 2}$, so $\mc{O}$ is a left-moving chiral or anti-chiral primary. Together with $(\wt h,\wt q)=({1\over 2},1)$ this gives precisely the $(c,c)$ and $(a,c)$ primaries:
\ie\label{defmatter}
e^{-{\wt\phi}}\,{\cal O}^{\rm NS}_{(\pm 1,1)}\,\lambda^{A'} \qquad (h,\wt h)=\Big({1\over 2},{1\over 2}\Big), \quad (q,\wt q)=(\pm 1,1),
\fe
as in (\ref{nsrmasslfrbsf}). Turning on such a deformation couples the internal SCFT to the left-moving gauge sector and generically breaks the holomorphic ${\cal N}=2$ structure, so it is a natural candidate to flow to a chiral $(0,2)$ theory.

\item {\bf No gauge fermions ($\Lambda=\mathbf{1}$): purely internal directions.}\\
Here $h(\mc{O})=1$ and the GSO condition forces $q(\mc{O})$ to be even. The unitarity bound implies $|q(\mc{O})|\le 2$, so $q(\mc{O})\in\{0,\pm 2\}$. This yields:
\begin{enumerate}
\item {\bf Spectral-flow images ($q=\pm 2$):} if $\mc{O}$ saturates the BPS bound at $h=1$ it is a chiral/anti-chiral primary with $q=\pm 2$, giving
\ie\label{defq2}
e^{-{\wt\phi}}\,{\cal O}'^{\rm NS}_{(\pm 2,1)} \qquad (h,\wt h)=(1,{1\over 2}), \quad (q,\wt q)=(\pm 2,1),
\fe
as in (\ref{nsrmasslfrbsf}).
\item {\bf Standard-embedding moduli ($q=0$ descendants):} the usual $(2,2)$-preserving marginal operators are the top components of short multiplets built on $(c,c)$ and $(a,c)$ primaries,
\ie\label{defmoduli}
e^{-{\wt\phi}}\,G_{-{1\over 2}}^\pm\,{\cal O}^{{\rm NS},\mp}_{(\mp 1,1)} \qquad (h,\wt h)=(1,{1\over 2}), \quad (q,\wt q)=(0,1),
\fe
as in (\ref{nsrmasslfrbsfneutral}).
\item {\bf Additional $q=0$ primaries (long multiplets):} finally, the internal SCFT may contain further operators with $(h,q)=(1,0)$ and $(\wt h,\wt q)=({1\over 2},1)$ not of the descendant form above (e.g.\ operators involving extra left-moving currents). If present, these define additional supersymmetric $(0,2)$ deformations.
\ie\label{defq0primary}
e^{-{\wt\phi}}\,{\cal P}^{\rm NS}_{(0,1)} \qquad (h,\wt h)=\Big(1,{1\over 2}\Big), \quad (q,\wt q)=(0,1),
\fe
\end{enumerate}

\item {\bf Two gauge fermions ($\Lambda=J^{A'B'}$): current-type directions.}\\
Here $h(\mc{O})=0$, so by unitarity $q(\mc{O})=0$ and $\mc{O}$ lies in the left vacuum module. Therefore such deformations exist if and only if the internal theory contains right-moving chiral primaries with $(\wt h,\wt q)=({1\over 2},1)$ in the left vacuum module, i.e.\ purely right-moving operators. When present, the corresponding $(-1)$ picture vertices take the form
\ie\label{defcurrent}
e^{-{\wt\phi}}\,J^{A'B'}\,\mc{U}_{(\wt h,\wt q)=({1\over 2},1)}(\bz),
\fe
and they generically break the holomorphic ${\cal N}=2$ structure while preserving the anti-holomorphic one.
\end{enumerate}

The CPT conjugates correspond to $(\wt h,\wt q)=({1\over 2},-1)$ together with exchanging $\wt G^-\leftrightarrow \wt G^+$ in (\ref{rFterm}); we will keep both $\wt q=\pm 1$ sectors explicit when discussing NS--NS--NS correlators, since picture-changing forces one to mix them.

\subsection{Yukawa couplings and cubic obstructions}

Given a set of candidate deformations classified in section 3.1, a necessary first check for a flat direction is F-flatness with respect to the cubic superpotential. On the worldsheet this is controlled by sphere 3-point functions with one NS vertex operator and two R vertex operators, as reviewed in section 2. We now analyze when the corresponding Yukawa couplings vanish purely by general selection rules of a standard embedding $(2,2)$ point.

\subsubsection*{General form of the Yukawa couplings}

For each supersymmetric marginal direction, we may choose unintegrated vertices (at zero momentum) of the form
\ie\label{defverts}
V_i^{\rm NS} = e^{-{\wt\phi}}\,\Lambda_i\,\mc{O}_i^{\rm NS}, \qquad
V_i^{\rm R} = e^{-{\wt\phi\over 2}}\,\wt S_{\A}\,\Lambda_i\,\mc{O}_i^{\rm R},
\fe
where $\Lambda_i$ is an operator in the free left-moving $\lambda^{A'}$ system (cf.\ the completeness result of section 3.1),
\ie
\Lambda_i \in \big\{ \mathbf{1},~\lambda^{A'},~J^{A'B'}\equiv :\lambda^{A'}\lambda^{B'}:\big\},
\fe
and $\mc{O}_i^{\rm NS}$ is an internal operator with right-moving quantum numbers $(\wt h,\wt q)=({1\over 2},1)$, while $\mc{O}_i^{\rm R}$ is its (NS,R) partner obtained by half-unit anti-holomorphic spectral flow, with $(\wt h,\wt q)=({3\over 8},-{1\over 2})$. The cubic Yukawa couplings are then extracted from correlators of the schematic form
\ie\label{yukawa3pt}
Y_{ijk} \;\propto\; \left\langle V_i^{\rm NS}(\infty)\,V_j^{\rm R}(1)\,V_k^{\rm R}(0)\right\rangle.
\fe
At a standard embedding point the $\lambda^{A'}$ system factorizes, so (\ref{yukawa3pt}) splits into a gauge-fermion correlator involving $\Lambda_i$ and an internal correlator involving $\mc{O}_i$.

\subsubsection*{Selection rule 1: free gauge-fermion Wick contractions}

By Wick's theorem, any correlator of the free fermions vanishes unless the total number of $\lambda^{A'}$ insertions is even. Moreover, since $J^{A'B'}$ is normal ordered, a single $J^{A'B'}$ insertion must contract with \emph{external} $\lambda$'s; in particular
\ie
\langle \lambda^{A'}\rangle = 0, \qquad \langle J^{A'B'}\rangle=0, \qquad
\langle \lambda^{A'}\lambda^{B'}\lambda^{C'}\rangle=0,
\fe
and $\langle J^{A'B'}\,{\cal X}\rangle$ vanishes unless ${\cal X}$ contains at least two additional $\lambda$ insertions. Concretely, the only potentially nonzero gauge-fermion structures at cubic order are of the schematic types
\ie\label{gaugecontractions}
\langle \lambda\,\lambda\rangle \neq 0,\qquad
\langle J\,\lambda\,\lambda\rangle \neq 0,\qquad
\langle J\,J\rangle \neq 0,\qquad
\langle J\,J\,J\rangle \neq 0,
\fe
with the obvious index contractions inherited from the free-fermion realization of the gauge current algebra.

Applied to the deformation multiplets of section 3.1, this immediately implies for example:
\begin{enumerate}
\item any Yukawa coupling involving \emph{three} $F=1$ multiplets (three single-$\lambda$ insertions) vanishes,
\item any Yukawa coupling with one current-type multiplet ($\Lambda=J$) and two $F=0$ multiplets ($\Lambda=\mathbf{1}$) vanishes,
\item a single current insertion must be accompanied by two additional $\lambda$ insertions (giving $\langle J\,\lambda\,\lambda\rangle$), while two or three current insertions can also contract among themselves (giving $\langle J\,J\rangle$ or $\langle J\,J\,J\rangle$).
\end{enumerate}

\subsubsection*{Selection rule 2: holomorphic $U(1)_R$ charge}

The holomorphic $U(1)_R$ current of the internal $(2,2)$ SCFT is conserved on the sphere, so the internal piece of (\ref{yukawa3pt}) can be nonzero only if the holomorphic charges satisfy
\ie\label{qsum0}
q_i + q_j + q_k = 0,
\fe
where $q_i$ is the holomorphic $U(1)_R$ charge of $\mc{O}_i$ (equivalently the left-moving charge of the corresponding chiral multiplet). For the supersymmetric marginal directions of section 3.1 one has $q\in\{ \pm 1, \pm 2, 0\}$, with $q=\pm 1$ precisely for the single-$\lambda$ matter multiplets, and $q=\pm 2,0$ for the $F=0$ multiplets.

Combining (\ref{qsum0}) with the gauge-fermion selection rule, we find that the only \emph{a priori} allowed Yukawa couplings involving the deformation multiplets are the following:
\begin{enumerate}
\item {\bf Two matter insertions:} $\langle \lambda\,\lambda\rangle\neq 0$ allows Yukawas with two $F=1$ multiplets and one $F=0$ multiplet. The charge constraint (\ref{qsum0}) leaves the possibilities
\ie
(q,q,q) \;=\; (+1,+1,-2),\quad (-1,-1,+2),\quad (+1,-1,0),
\fe
i.e.\ the chiral-ring Yukawas discussed in section 2, together with the $q=0$ case.
\item {\bf One current insertion:} $\langle J\,\lambda\,\lambda\rangle\neq 0$ forces two $F=1$ multiplets and one current-type multiplet ($\Lambda=J$). Then (\ref{qsum0}) requires the two matter multiplets to have opposite charges $q=\pm 1$.
\item {\bf No matter insertions:} if all three multiplets have $\Lambda=\mathbf{1}$, the gauge-fermion correlator is trivial and only the holomorphic charge condition remains. Since $q\in\{0,\pm 2\}$ in this sector, (\ref{qsum0}) permits only
\ie
(q,q,q)\;=\;(0,0,0)\quad{\rm or}\quad (0,+2,-2).
\fe
\item {\bf Two or three current insertions:} $\langle J\,J\rangle\neq 0$ (respectively $\langle J\,J\,J\rangle\neq 0$) allows Yukawas with two (respectively three) current-type multiplets. Since current-type multiplets have $q=0$, the remaining multiplet (if any) must also have $q=0$ by (\ref{qsum0}).
\end{enumerate}
All other cubic couplings among the deformation multiplets vanish identically by the two selection rules above.

\subsubsection*{Ward identity: vanishing for standard-embedding moduli}

We now show that the $q=0$ \emph{standard-embedding moduli} (the $F=0$ descendants of the form $G_{-{1\over 2}}^\pm{\cal O}^{{\rm NS},\mp}$) have vanishing cubic Yukawa couplings with pairs of deformation multiplets whose holomorphic factors are chiral/anti-chiral primaries (in particular the $q=\pm 1$ and $q=\pm 2$ multiplets, as well as the current-type multiplets whose holomorphic factor is the identity). Consider the holomorphic part of a correlator of the form
\ie\label{Gdesc3pt}
\left\langle \big(G_{-{1\over 2}}^\pm \phi\big)(\infty)\,\psi(1)\,\chi(0)\right\rangle,
\fe
where $\phi$ is a chiral primary with $q(\phi)=\mp 1$ (so that $G^\pm_{-{1\over 2}}\phi$ has $q=0$), and $\psi,\chi$ are holomorphic primaries with charges satisfying $q(\psi)+q(\chi)=0$ as in (\ref{qsum0}). Representing the insertion at $z=\infty$ as a contour integral $\oint {dz\over 2\pi i}\,z\,G^\pm(z)$ and shrinking the contour, one obtains a sum of terms where $G^\pm$ acts on the operators at $z=1$ and $z=0$. With the choice of $f(z)=z$, the contribution at $z=0$ involves the positive mode $G^\pm_{1\over 2}$ acting on $\chi$, which vanishes because $\chi$ is a primary. The remaining term involves $G^\pm_{-{1\over 2}}$ acting on $\psi$.

For the Yukawa couplings between the BPS deformation multiplets, charge conservation forces $\psi$ and $\chi$ to be a chiral/anti-chiral pair of opposite holomorphic charges. In particular, one of $\psi,\chi$ is a chiral primary annihilated by $G^+_{-{1\over 2}}$, and one is an anti-chiral primary annihilated by $G^-_{-{1\over 2}}$. Choosing the sign in (\ref{Gdesc3pt}) to match the standard-embedding modulus, and placing the operator annihilated by $G^\pm_{-{1\over 2}}$ at $z=1$, the remaining term therefore vanishes. We conclude that the standard-embedding moduli have no cubic Yukawa couplings to pairs of BPS deformation multiplets. This is the generalization of the example worked out explicitly in section 2.

In contrast, $q=0$ \emph{primary} multiplets (when present) are not protected by this Ward identity, and their Yukawa couplings of type (3) above are generically nonzero unless forbidden by additional symmetries.

\subsubsection*{NS--NS--NS correlators and conformal perturbation theory}

The Yukawa correlators (\ref{yukawa3pt}) test cubic superpotential obstructions. Exact marginality, however, is governed by the OPE of the \emph{integrated} marginal operators that appear in the supersymmetric deformation (\ref{rFterm}), and the corresponding OPE coefficients are extracted from sphere 3-point functions of three NS insertions.

In the BRST description the NS--NS--NS correlator requires one insertion of the PCO, because on the sphere the total right-moving picture number must be $-2$. Choosing the picture assignment $(0,-1,-1)$, we write the chiral and anti-chiral $(-1)$-picture NS vertices as
\ie
V_i^{(-1)} = e^{-{\wt\phi}}\,\Lambda_i\,\mc{O}_i^{\rm NS},\qquad
\overline V_i^{(-1)} = e^{-{\wt\phi}}\,\Lambda_i\,\overline{\mc{O}}{}_i^{\rm NS},
\fe
where $\mc{O}_i^{\rm NS}$ has $(\wt h,\wt q)=({1\over 2},1)$ and $\overline{\mc{O}}{}_i^{\rm NS}$ has $(\wt h,\wt q)=({1\over 2},-1)$. The PCO is
\ie
X(\bz)\equiv \{Q_{\rm BRST},\wt\xi(\bz)\} \;=\; e^{{\wt\phi}(\bz)}\big(\wt G^+ + \wt G^-\big)(\bz) + \cdots,
\fe
so it contains an explicit insertion of the right-moving supercurrent. The 0-picture vertices are obtained by applying $X$:
\ie\label{V0fromPCO}
V_i^{(0)} \equiv X\cdot V_i^{(-1)} \;\sim\; \wt G^-_{-{1\over 2}}\big(\Lambda_i\,\mc{O}_i^{\rm NS}\big),\qquad
\overline V_i^{(0)} \equiv X\cdot \overline V_i^{(-1)} \;\sim\; \wt G^+_{-{1\over 2}}\big(\Lambda_i\,\overline{\mc{O}}{}_i^{\rm NS}\big),
\fe
where only the displayed term in $\wt G^\pm$ contributes because $\wt G^+_{-{1\over 2}}\mc{O}_i^{\rm NS}=0$ and $\wt G^-_{-{1\over 2}}\overline{\mc{O}}{}_i^{\rm NS}=0$ by right chirality. The conformal-perturbation-theory coefficients $C^i{}_{jk}$ of section~\ref{sec:cpt_obstructions} are then proportional to NS--NS--NS correlators of the schematic form
\ie\label{NSNSNS3pt}
C^i{}_{j\bar k}\;\propto\;\left\langle V_i^{(0)}(\infty)\,V_j^{(-1)}(1)\,\overline V_k^{(-1)}(0)\right\rangle,
\fe
together with permutations and the conjugate correlators with $V\leftrightarrow \overline V$. At a standard embedding $(2,2)$ point, conservation of the right-moving $U(1)_R$ current implies that the two $(-1)$-picture insertions must have opposite right charges,
\ie\label{qtildeSelection}
\wt q_j + \wt q_k = 0,
\fe
so one must be in the $\wt q=1$ sector and the other in the $\wt q=-1$ sector.

\subsection{Possibly non-vanishing 3-point functions}

The selection rules of section 3.2 leave only a small set of 3-point functions that are not forced to vanish at a generic standard embedding $(2,2)$ point. We distinguish (i) NS--R--R correlators that compute spacetime Yukawa couplings, and (ii) NS--NS--NS correlators (with one PCO insertion) that compute the OPE coefficients entering conformal perturbation theory. Throughout, ``possibly non-vanishing'' means not ruled out by the free-fermion contractions (\ref{gaugecontractions}) and holomorphic charge conservation (\ref{qsum0}); in addition, the Yukawas are subject to the Ward-identity vanishing for the $q=0$ descendant moduli, while the NS--NS--NS correlators are subject to the right-charge condition (\ref{qtildeSelection}).

\subsubsection*{NS--R--R Yukawas}

\begin{enumerate}
\item {\bf Chiral-ring Yukawas (two matter and one $q=\mp 2$ multiplet).}\\
These have the schematic form
\ie
\langle \lambda^{A'}\lambda^{B'}\rangle~\left\langle {\cal O}'^{{\rm NS},(\mp 2,1)}~{\cal O}^{{\rm R},(\pm 1,-{1\over 2})}~{\cal O}^{{\rm R},(\pm 1,-{1\over 2})} \right\rangle^{(2,2)~{\rm SCFT}},
\fe
corresponding to $(q,q,q)=(+1,+1,-2)$ and $(q,q,q)=(-1,-1,+2)$.

\item {\bf Matter--neutral couplings with $(q,q,q)=(+1,-1,0)$.}\\
If the neutral multiplet is a $q=0$ \emph{primary} (a long multiplet of section 3.1), then nothing forbids
\ie
\langle \lambda^{A'}\lambda^{B'}\rangle~\left\langle {\cal P}^{{\rm NS},(0,1)}~{\cal O}^{{\rm R},(+1,-{1\over 2})}~{\cal O}^{{\rm R},(-1,-{1\over 2})} \right\rangle^{(2,2)~{\rm SCFT}},
\fe
where ${\cal P}^{{\rm NS},(0,1)}$ denotes an internal operator with $(h,q)=(1,0)$ and $(\wt h,\wt q)=({1\over 2},1)$, and ${\cal P}^{{\rm R},(0,-{1\over 2})}$ is its (NS,R) partner. By contrast, if the $q=0$ insertion is a standard-embedding modulus of descendant form $G_{-{1\over 2}}^\pm{\cal O}^{{\rm NS},\mp}$, the corresponding Yukawa coupling vanishes by the Ward identity of section 3.2.

\item {\bf Purely internal $F=0$ Yukawas.}\\
In the $\Lambda=\mathbf{1}$ sector, charge conservation allows only $(q,q,q)=(0,0,0)$ and $(0,+2,-2)$. For the $(0,+2,-2)$ pattern, the Ward identity of section 3.2 implies vanishing if the $q=0$ insertion is a standard-embedding modulus descendant, so a nonzero coupling requires a $q=0$ primary. By contrast, the $(0,0,0)$ couplings are not constrained by that Ward identity and may involve any combination of $q=0$ multiplets (descendants and/or primaries), depending on the model. Examples of internal correlators appearing are
\ie
\left\langle {\cal P}^{{\rm NS},(0,1)}~{\cal O}'^{{\rm R},(+2,-{1\over 2})}~{\cal O}'^{{\rm R},(-2,-{1\over 2})} \right\rangle^{(2,2)~{\rm SCFT}},
\qquad
\left\langle {\cal P}^{{\rm NS},(0,1)}~{\cal P}^{{\rm R},(0,-{1\over 2})}~{\cal P}^{{\rm R},(0,-{1\over 2})} \right\rangle^{(2,2)~{\rm SCFT}},
\fe
together with permutations.

\item {\bf Yukawas involving current-type multiplets.}\\
Current-type multiplets have $\Lambda=J^{A'B'}$ and holomorphic identity in the internal factor (section 3.1). There are then two qualitatively different possibilities:
\begin{enumerate}
\item {\bf One current and two matter:} a potentially nonzero correlator is
\ie
\langle J^{A'B'}\,\lambda^{C'}\,\lambda^{D'}\rangle~
\left\langle \mc{U}^{{\rm NS},(0,1)}~{\cal O}^{{\rm R},(+1,-{1\over 2})}~{\cal O}^{{\rm R},(-1,-{1\over 2})} \right\rangle^{(2,2)~{\rm SCFT}},
\fe
with the matter multiplets carrying opposite charges $q=\pm 1$.
\item {\bf Three currents:} with three current-type multiplets one has
\ie
\langle J\,J\,J\rangle~\left\langle \mc{U}_i^{{\rm NS},(0,1)}~\mc{U}_j^{{\rm R},(0,-{1\over 2})}~\mc{U}_k^{{\rm R},(0,-{1\over 2})} \right\rangle^{(2,2)~{\rm SCFT}},
\fe
which is model-dependent but not forbidden by the general selection rules.
\end{enumerate}
A Yukawa with exactly two current-type multiplets has a nonzero gauge-fermion factor $\langle J\,J\rangle$, but since the current-type insertions have holomorphic identity in the internal factor, the internal holomorphic piece reduces to a one-point function of the remaining insertion. This vanishes on the sphere unless that insertion is the identity, which is not the case for a marginal deformation.
\end{enumerate}

\subsubsection*{NS--NS--NS correlators (CPT/OPE coefficients)}

The NS--NS--NS correlators relevant for conformal perturbation theory are of the form (\ref{NSNSNS3pt}), i.e.\ they involve a 0-picture vertex obtained by picture-changing and thus contain an explicit right-moving supercurrent insertion. The remaining possibilities consistent with (\ref{gaugecontractions}), (\ref{qsum0}), and (\ref{qtildeSelection}) may be organized as follows (together with permutations and conjugates):
\begin{enumerate}
\item {\bf Two matter and one $F=0$ insertion ($\langle \lambda\,\lambda\rangle\neq 0$).}\\
The holomorphic charge patterns are the same as for the Yukawas,
\ie
(q,q,q)=(+1,+1,-2),\quad (-1,-1,+2),\quad (+1,-1,0),
\fe
	but now with one of the $(-1)$-picture matter insertions replaced by its right-CPT conjugate to satisfy (\ref{qtildeSelection}). A representative correlator is
	\ie
	\langle \lambda^{A'}\lambda^{B'}\rangle~\left\langle \big(\wt G^-_{-{1\over 2}}\mc{O}_i^{\rm NS}\big)~{\cal O}_j^{\rm NS}~\overline{\cal O}{}_k^{\rm NS} \right\rangle^{(2,2)~{\rm SCFT}},
	\fe
	where ${\cal O}_{j,k}^{\rm NS}$ are matter operators (\ref{defmatter}) and the $F=0$ insertion is taken from (\ref{defq2}), (\ref{defmoduli}), or (\ref{defq0primary}) according to the desired $q$.
	In the $(q,q,q)=(+1,-1,0)$ channel, if the $q=0$ insertion is a standard-embedding modulus descendant of the form (\ref{defmoduli}), the holomorphic ${\cal N}=2$ Ward identity of section 3.2 forces the correlator to vanish; thus a nonzero coefficient requires the $q=0$ insertion to be a $q=0$ primary (\ref{defq0primary}).

\item {\bf Purely internal $F=0$ correlators ($\Lambda=\mathbf{1}$).}\\
	Here the gauge-fermion factor is trivial and holomorphic charge conservation restricts to $(q,q,q)=(0,0,0)$ and $(0,+2,-2)$ (with the $q=0$ sector including both descendants (\ref{defmoduli}) and possible primaries (\ref{defq0primary})). Representative internal correlators are
	\ie
	\left\langle \big(\wt G^-_{-{1\over 2}}{\cal P}^{\rm NS}_{(0,1)}\big)~{\cal P}^{\rm NS}_{(0,1)}~\overline{\cal P}{}^{\rm NS}_{(0,-1)} \right\rangle^{(2,2)~{\rm SCFT}},
	\qquad
	\left\langle \big(\wt G^-_{-{1\over 2}}{\cal O}'^{\rm NS}_{(+2,1)}\big)~{\cal O}'^{\rm NS}_{(-2,1)}~\overline{\cal P}{}^{\rm NS}_{(0,-1)} \right\rangle^{(2,2)~{\rm SCFT}},
	\fe
	together with permutations and the variants where the $(-1)$-picture insertion with $\wt q=-1$ is one of the $q=\pm 2$ operators instead.
	As in the Yukawa discussion, the holomorphic ${\cal N}=2$ Ward identity implies that in the $(q,q,q)=(0,+2,-2)$ channel the correlator vanishes if the $q=0$ insertion is a standard-embedding modulus descendant (\ref{defmoduli}); a nonzero coefficient therefore requires a $q=0$ primary (\ref{defq0primary}). No such Ward-identity protection applies to the $(0,0,0)$ family.

\item {\bf One current and two matter ($\langle J\,\lambda\,\lambda\rangle\neq 0$).}\\
With one current-type multiplet and two matter multiplets, holomorphic charge conservation forces the matter charges to be opposite, $q=\pm 1$. A representative correlator is
\ie
\langle J^{A'B'}\,\lambda^{C'}\,\lambda^{D'}\rangle~
\left\langle \big(\wt G^-_{-{1\over 2}}\mc{U}^{\rm NS}\big)~{\cal O}^{\rm NS}_{(+1,1)}~\overline{\cal O}{}^{\rm NS}_{(-1,-1)} \right\rangle^{(2,2)~{\rm SCFT}},
\fe
with the understanding that one of the two matter insertions is chosen from the $\wt q=-1$ sector to satisfy (\ref{qtildeSelection}).

\item {\bf Three currents ($\langle J\,J\,J\rangle\neq 0$).}\\
With three current-type multiplets, holomorphic charge conservation is automatic (all $q=0$) and the only model-independent constraint beyond gauge contractions is (\ref{qtildeSelection}). A representative correlator is
\ie
\langle J\,J\,J\rangle~\left\langle \big(\wt G^-_{-{1\over 2}}\mc{U}_i^{\rm NS}\big)~\mc{U}_j^{\rm NS}~\overline{\mc{U}}{}_k^{\rm NS} \right\rangle^{(2,2)~{\rm SCFT}},
\fe
where $\mc{U}^{\rm NS}$ denotes a purely right-moving chiral primary in the left vacuum module (section 3.1), and $\overline{\mc{U}}{}^{\rm NS}$ is its right-CPT conjugate.
\end{enumerate}
As for the Yukawas, a correlator with exactly two current-type multiplets has a nonzero gauge-fermion factor $\langle J\,J\rangle$, but the internal holomorphic piece reduces to a one-point function of the remaining insertion and therefore vanishes for a marginal deformation.

The vanishing (or not) of the Yukawa correlators above provides the first, model-independent obstruction test for deformations leaving the standard embedding locus. Independently, the NS--NS--NS correlators determine the leading (quadratic) beta functions in conformal perturbation theory; higher-order obstructions are then governed by higher-point functions, which we organize in section~\ref{sec:cpt_obstructions}.

\subsection{Large-volume interpretation}

When the internal $(2,2)$ SCFT admits a weakly coupled large-volume description as a supersymmetric nonlinear sigma model on a Calabi--Yau threefold $X$ with background gauge fields \cite{Hull:1985jv,Ellwanger:1988cc}, the standard embedding point corresponds to choosing the left-moving gauge bundle to coincide with the tangent bundle,
\ie
E \;=\; T_X,
\fe
so that the worldsheet theory has an enhanced $(2,2)$ structure \cite{Sen:1986mg}. In this description the right-moving chiral primary of quantum numbers $(\wt h,\wt q)=({1\over 2},1)$ that appears in the general supersymmetric deformation (\ref{rFterm}) is represented schematically by a right-moving fermion $\wt\psi^{\bar i}(\bz)$ (with $i,\bar i$ target-space indices), and its anti-holomorphic descendant satisfies
\ie\label{GpsiX}
\wt G^-_{-{1\over 2}}\,\wt\psi^{\bar i} \;\sim\; \bar\partial X^{\bar i} + \cdots.
\fe
These free-field representatives underlie both types of three-point functions discussed above: the NS--R--R Yukawas and the NS--NS--NS correlators relevant for conformal perturbation theory (with one PCO insertion). We therefore present the large-volume operator dictionary once and then focus on the extra ingredient for CPT, namely the picture-changing replacement (\ref{GpsiX}).
The geometric meaning of the deformation classes of section 3.1 follows by inserting explicit free-field representatives of the internal chiral primaries. A useful general rule is that a Dolbeault class in
\ie
H^{1}(X,V)\;\equiv\; H^{0,1}_{\bar\partial}(X,V)
\;=\; {\ker(\bar\partial:\Omega^{0,1}(V)\to \Omega^{0,2}(V))\over \mathrm{im}(\bar\partial:\Omega^{0,0}(V)\to \Omega^{0,1}(V))}
\fe
may be represented by a $\bar\partial$-closed $(0,1)$-form $\alpha_{\bar i}(X)$ valued in a bundle $V$ (more generally, $\bar\partial$ is the Dolbeault operator $\bar\partial_V$ defining the holomorphic structure of $V$). The corresponding worldsheet operator is obtained by contracting $\alpha_{\bar i}$ with a right-moving fermion $\wt\psi^{\bar i}$ (and, when needed to reach left-moving weight $h=1$, with appropriate left-moving fermions). The freedom to shift $\alpha_{\bar i}\to \alpha_{\bar i}+\bar\partial_{\bar i}\beta$ corresponds to adding a BRST-exact operator, so the vertex depends only on the cohomology class.

\begin{enumerate}
\item {\bf $(2,2)$-preserving moduli ($\Lambda=\mathbf{1}$, $q=0$ descendants).}\\
The familiar K\"ahler and complex-structure moduli correspond to the $q=0$ descendant multiplets of section 3.1, whose unintegrated $(-1)$ picture vertices are of the form (\ref{defmoduli}). In the free-field/large-volume limit, the underlying $(c,c)$ or $(a,c)$ primaries with $(h,\wt h)=({1\over 2},{1\over 2})$ and $(q,\wt q)=(\pm 1,1)$ are represented by operators built from one left-moving and one right-moving fermion, e.g.\ $J_{i\bar j}\,\psi^i(z)\wt\psi^{\bar j}(\bz)$ for a K\"ahler deformation. Acting with $G^\pm_{-{1\over 2}}$ produces the marginal $(1,1)$ operators that deform the metric and $B$-field (and similarly for complex structure).
\begin{enumerate}
\item {\bf K\"ahler moduli:} a harmonic representative $\omega_{i\bar j}\in H^{1,1}(X)\simeq H^{1}(X,T_X^\vee)$ gives the $(c,c)$ primary
\ie
{\cal O}^{{\rm NS},(+1,1)} \;\sim\; \omega_{i\bar j}(X)\,\psi^i(z)\wt\psi^{\bar j}(\bz),
\fe
whose descendant in (\ref{defmoduli}) produces the usual supersymmetric deformation of $g_{i\bar j}$ and $B_{i\bar j}$,
\ie
\delta S_{g+B}\;\sim\;\int d^2z~\omega_{i\bar j}(X)\,\partial X^i\,\bar\partial X^{\bar j} + \cdots.
\fe
\item {\bf Complex-structure moduli:} a Beltrami differential $\mu_{\bar i}{}^{j}\in H^{1}(X,T_X)\simeq H^{2,1}(X)$ gives the $(a,c)$ primary
\ie
{\cal O}^{{\rm NS},(-1,1)} \;\sim\; \mu_{\bar i}{}^{j}(X)\,\psi_{j}(z)\wt\psi^{\bar i}(\bz),
\fe
	where $\psi_{j}=g_{j\bar k}\psi^{\bar k}$. Acting with $G^\pm_{-{1\over 2}}$ and $\wt G^-_{-{1\over 2}}$ as in (\ref{defmoduli}) produces the corresponding marginal deformation of the sigma-model action (schematically $\delta S\sim\int \mu_{\bar i}{}^{j}\,\partial X^{j}\bar\partial X^{\bar i}+\cdots$), implementing the complex-structure variation.
	At leading order the picture-changed 0-picture insertion appearing in the NS--NS--NS correlator (\ref{NSNSNS3pt}) is precisely this $(1,1)$ operator, obtained by replacing $\wt\psi^{\bar i}\to \bar\partial X^{\bar i}$ via (\ref{GpsiX}). The corresponding quadratic CPT obstruction is represented by the Kodaira--Spencer bracket of Beltrami differentials,
	\ie
	\relax [\mu_1,\mu_2]_{\bar i\bar j}{}^{k}
	\;\equiv\;
	\mu_{1[\bar i}{}^{\ell}\,\partial_{\ell} \mu_{2\bar j]}{}^{k}
	-
	\mu_{2[\bar i}{}^{\ell}\,\partial_{\ell} \mu_{1\bar j]}{}^{k}
	\;\in\;\Omega^{0,2}(X,T_X),
	\fe
	whose cohomology class in $H^{2}(X,T_X)$ is the familiar integrability obstruction.
	\end{enumerate}
	
	\item {\bf Bundle deformations (current-type multiplets with $\Lambda=J$).}\\
	Deforming away from the standard embedding while preserving $(0,2)$ supersymmetry means deforming the holomorphic structure of the left-moving bundle $E$. Infinitesimally this is described by $\delta A\in\Omega^{0,1}(\mathrm{End}\,E)$ modulo gauge transformations, with linearized holomorphy condition $\bar\partial_E \delta A=0$ and equivalence $\delta A\sim \delta A+\bar\partial_E\epsilon$, hence a class in $H^{1}(X,\mathrm{End}\,E)$. On the worldsheet, such deformations appear precisely as the current-type multiplets of section 3.1, with unintegrated vertices of the form (\ref{defcurrent}): the left-moving factor is a gauge current $J^{A'B'}=: \lambda^{A'}\lambda^{B'}:$, while the right-moving chiral primary in the left vacuum module is represented at large volume by $\mc{U}(\bz)=\delta A_{\bar i}^{A'B'}(X)\,\wt\psi^{\bar i}(\bz)$. Using (\ref{GpsiX}), the integrated perturbation (\ref{rFterm}) then reduces to the familiar NLSM deformation of the gauge connection in the left-moving kinetic term,
	\ie
	\delta S_E \;\sim\; \int d^2z~\delta A_{\bar i}^{A'B'}(X)\,\bar\partial X^{\bar i}\,J^{A'B'}(z) + \cdots.
		\fe
		Equivalently, in the NS--NS--NS correlator (\ref{NSNSNS3pt}) the required 0-picture insertion is obtained by picture-changing the $(-1)$-picture vertex $e^{-{\wt\phi}}J^{A'B'}\delta A_{\bar i}^{A'B'}\wt\psi^{\bar i}$, which again implements $\wt\psi^{\bar i}\to \bar\partial X^{\bar i}$ at leading order. The corresponding quadratic CPT obstructions are encoded in the bundle bracket
		\ie
		\relax [\delta A_1,\delta A_2]_{\bar i\bar j}
	\;\equiv\;
	[\delta A_{1\bar i},\delta A_{2\bar j}]-
	[\delta A_{1\bar j},\delta A_{2\bar i}]
		\;\in\;\Omega^{0,2}(X,\mathrm{End}\,E),
		\fe
		which is the large-volume avatar of the structure constant extracted from the NS--NS--NS sphere correlator (\ref{NSNSNS3pt}), e.g.
		\[
		\left\langle V_{\delta A_3}^{(0)}(\infty)\,V_{\delta A_1}^{(-1)}(1)\,\overline V_{\delta A_2}^{(-1)}(0)\right\rangle
		\]
		for current-type vertices $V_{\delta A}^{(-1)}\sim e^{-{\wt\phi}}J^{A'B'}\delta A_{\bar i}^{A'B'}\,\wt\psi^{\bar i}$ and $V^{(0)}=X\cdot V^{(-1)}$ as in (\ref{V0fromPCO}).
		and in the mixed curvature contraction with a complex-structure deformation,
		\ie
		(\iota_{\mu}{\cal F})_{\bar i\bar j}
	\;\equiv\;
	\mu_{[\bar i}{}^{k}\,{\cal F}_{\bar j]k}
		\;\in\;\Omega^{0,2}(X,\mathrm{End}\,E),
			\fe
			which for a generic holomorphic bundle background is obtained by projecting the mixed NS--NS--NS correlator (\ref{NSNSNS3pt}) onto a bundle direction, e.g.
			\[
			\left\langle V_{\delta A_3}^{(0)}(\infty)\,V_{\mu}^{(-1)}(1)\,\overline V_{\delta A_2}^{(-1)}(0)\right\rangle
			\]
			with $V_{\mu}^{(-1)}\sim e^{-{\wt\phi}}\,G^+_{-{1\over 2}}\big(\mu_{\bar i}{}^{j}\psi_{j}\,\wt\psi^{\bar i}\big)\sim e^{-{\wt\phi}}\,\mu_{\bar i}{}^{j}\,\partial X^{j}\,\wt\psi^{\bar i}+\cdots$ the complex-structure modulus vertex of (\ref{defmoduli}).
			At the strict standard embedding $(2,2)$ point, the gauge-fermion factor $\langle J\,J\rangle$ is nonzero but the internal \emph{holomorphic} factor reduces to a one-point function, so the full correlator vanishes by the holomorphic one-point selection rule (section 3.3), in accord with $\alpha_{E_0}=0$; the first nontrivial mixed lifting induced by turning on a small current-type background is therefore captured by the corresponding integrated four-point functions, cf.\ section \ref{sec:4pt_selection}. The associated cohomology class is the Atiyah map (\ref{eq:Atiyah_map}) controlling which complex-structure directions survive away from the standard embedding locus.
		
		\item {\bf Charged Higgs-branch directions (single-$\lambda$ matter multiplets).}\\
	The deformations with $\Lambda=\lambda^{A'}$ (cf.\ (\ref{defmatter})) correspond to turning on backgrounds for spacetime chiral multiplets that are charged under the unbroken 4D gauge symmetry at the standard embedding point. In the large-volume sigma-model language these are bundle-valued ``matter'' fluctuations. For a general heterotic compactification with bundle $E$, massless charged matter in a representation $R$ is counted by $H^{1}(X,V_R)$, where $V_R$ is the associated bundle (the conditions $\bar\partial_{V_R}\alpha=0$ and $\alpha\sim \alpha+\bar\partial_{V_R}\beta$ are the linearized equations of motion and gauge equivalences). At the level of free fields, a representative $\alpha_{\bar i}(X)\in H^{1}(X,V_R)$ leads to a vertex of the schematic form
\ie
V_{\rm matter}^{\rm NS}\;\sim\; e^{-{\wt\phi}}\,\lambda\,\alpha_{\bar i}(X)\,\psi(z)\,\wt\psi^{\bar i}(\bz),
\fe
with the additional left-moving fermion $\psi$ supplied by the internal $(2,2)$ SCFT so that the total left weight is $h=1$. In particular, at the standard embedding $E=T_X$ the relevant associated bundles are $T_X$ and $T_X^\vee$, giving cohomology groups $H^{1}(X,T_X)$ and $H^{1}(X,T_X^\vee)$, and one may take the free-field representatives
\ie
\lambda^{A'}\,\mu_{\bar i}{}^{j}(X)\,\psi_{j}\,\wt\psi^{\bar i}, \qquad
\lambda^{A'}\,\omega_{i\bar j}(X)\,\psi^{i}\,\wt\psi^{\bar j},
\fe
	for the $q=-1$ and $q=+1$ matter multiplets respectively. Applying $\wt G^-_{-{1\over 2}}$ as in (\ref{rFterm}) and using (\ref{GpsiX}) gives the corresponding marginal perturbation schematically as
	\ie
	\delta S_{\rm matter}\;\sim\;
	\int d^2z~\lambda^{A'}\,\mu_{\bar i}{}^{j}(X)\,\psi_{j}(z)\,\bar\partial X^{\bar i}
	\qquad{\rm or}\qquad
	\int d^2z~\lambda^{A'}\,\omega_{i\bar j}(X)\,\psi^{i}(z)\,\bar\partial X^{\bar j}
	\;+\;\cdots.
	\fe
	In either case the integrand has $(h,\wt h)=(1,1)$: $\lambda^{A'}$ and $\psi$ each contribute $h={1\over 2}$ while $\bar\partial X^{\bar i}$ contributes $\wt h=1$.
	Giving vevs to such fields generically moves onto a Higgs branch and can be viewed geometrically as deforming the bundle away from $T_X$ by turning on off-diagonal components of the gauge connection, thereby breaking the holomorphic $(2,2)$ structure to $(0,2)$.
\end{enumerate}

Finally, section 3.1 also allows two further types of supersymmetric $(0,2)$ deformations, whose existence and interpretation are more model-dependent:
\begin{enumerate}
\item {\bf Spectral-flow BPS multiplets ($q=\pm 2$, $\Lambda=\mathbf{1}$).}\\
These are the chiral/anti-chiral primaries of the form (\ref{defq2}). In a large-volume sigma model they may be represented by operators with two left-moving fermions and one right-moving fermion. Concretely, one may take
\ie
{\cal O}'^{{\rm NS},(+2,1)} \;\sim\; \rho_{\bar i\,jk}(X)\,\psi^j(z)\psi^k(z)\,\wt\psi^{\bar i}(\bz), \qquad
{\cal O}'^{{\rm NS},(-2,1)} \;\sim\; \rho_{\bar i}{}^{jk}(X)\,\psi_{j}(z)\psi_{k}(z)\,\wt\psi^{\bar i}(\bz),
\fe
with $\rho_{\bar i\,jk}\in H^{1}(X,\wedge^2 T_X^\vee)$ and $\rho_{\bar i}{}^{jk}\in H^{1}(X,\wedge^2 T_X)$. In a general heterotic compactification with bundle $E$, analogous operators are naturally valued in $\wedge^2 E^\vee$ and $\wedge^2 E$.
\;Their integrated F-terms follow by applying $\wt G^-_{-{1\over 2}}$ and using (\ref{GpsiX}), giving schematically
\ie
\delta S_{q=\pm 2}\;\sim\;\int d^2z~\rho_{\bar i}(\cdots)\,\psi\psi\,\bar\partial X^{\bar i} + \cdots.
\fe
\item {\bf Neutral primaries ($q=0$, $\Lambda=\mathbf{1}$).}\\
These are the long-multiplet directions of the form (\ref{defq0primary}). In a geometric sigma model they arise only when there are additional holomorphic weight-one primaries of zero $U(1)_R$ charge (beyond those in the ${\cal N}=2$ algebra). At free-field points (e.g.\ toroidal orbifolds) one can take such a primary to be a holomorphic current $j_a(z)$ (for instance $j_a=i\partial Y^a$ on a torus), and then ${\cal P}^{\rm NS}_{(0,1)}\sim j_a(z)\,\wt\psi^{\bar i}(\bz)$ describes a current--current marginal deformation, including Wilson-line moduli.
\end{enumerate}

\section{Large-volume constraints near the standard embedding}\label{sec:geom_constraints}

In this section we assume the standard-embedding point admits a weakly coupled large-volume description as a heterotic NLSM on a smooth Calabi--Yau threefold $X$ with vanishing tree-level flux and gauge bundle
\ie
E \;=\; T_X,
\fe
so that the worldsheet theory has the accidental holomorphic ${\cal N}=2$ symmetry of a $(2,2)$ model. We then ask what happens to the familiar $(2,2)$ moduli when we turn on a \emph{small} supersymmetric $(0,2)$ deformation away from $E=T_X$: which directions are obstructed already at first order in the large-volume expansion, what geometric conditions must hold, and how can those conditions be recognized in the abstract CFT language of section 3.

\subsection{Standard embedding at large volume: the $(Z,Y,\Lambda)$ description}

At tree level in $\alpha'$, a convenient worldsheet description of gauge-neutral first-order deformations of a heterotic NLSM is given by Melnikov and Sharpe \cite{Melnikov:2011aa}. On a Calabi--Yau target with $H=0$, the relevant right-chiral operator is represented classically by a triple
\ie\label{eq:geom_triple}
(Z,Y,\Lambda)\in \Omega^{0,1}(X,T_X)\oplus \Omega^{1,1}(X)\oplus \Omega^{0,1}(X,\mathrm{End}\,E),
\fe
defined modulo equivalences, and chirality (up to equations of motion) imposes coupled conditions that reduce at the standard embedding point $E=T_X$ to \cite{Melnikov:2011aa}
\ie\label{eq:geom_cond}
\bar\partial Z \;=\; 0,\qquad
\bar\partial Y \;=\; 0,\qquad
\bar\partial \Lambda \;=\; \iota_Z R,
\fe
where $R$ is the $(1,1)$ curvature of the K\"ahler connection on $T_X$. On the $(2,2)$ locus there is a further simplification: defining
\ie\label{eq:Lambda_tilde}
\widetilde\Lambda \;\equiv\; \Lambda - \nabla Z,
\fe
one finds $\bar\partial\,\widetilde\Lambda=0$ (and $\widetilde\Lambda\sim \widetilde\Lambda+\bar\partial\,\widetilde\lambda$), so the first-order deformations split canonically as
\ie
\relax [Z]\in H^1(X,T_X),\qquad [Y]\in H^1(X,T_X^\vee),\qquad [\widetilde\Lambda]\in H^1(X,\mathrm{End}\,T_X).
\fe
In the operator classification of section 3.1, the $(Z,Y)$ directions correspond to the $q=0$ descendant moduli (\ref{defmoduli}), while the $\widetilde\Lambda$ directions correspond to the current-type multiplets (\ref{defcurrent}) and are precisely the gauge-neutral deformations that leave the standard embedding locus already at first order.

\subsection{Infinitesimal departure: the linearized holomorphy/Atiyah constraint}

Now turn on a small $(0,2)$ deformation away from the standard embedding by deforming the gauge bundle to a nearby holomorphic bundle $E_\varepsilon$ (as a smooth bundle, still isomorphic to $T_X$) with curvature ${\cal F}_\varepsilon$ and ${\cal F}_0=R$. The Melnikov--Sharpe chirality condition for a combined complex-structure deformation $[Z]\in H^1(X,T_X)$ and bundle deformation $\Lambda$ takes the familiar form \cite{Melnikov:2011aa}
\ie
\bar\partial_{E_\varepsilon}\Lambda \;=\; \iota_Z{\cal F}_\varepsilon.
\fe
For a fixed holomorphic bundle $E_\varepsilon$, this equation admits a solution $\Lambda$ if and only if the cohomology class
\ie\label{eq:Atiyah_map}
\alpha_{E_\varepsilon}([Z]) \;\equiv\; [\iota_Z{\cal F}_\varepsilon]\;\in\; H^2(X,\mathrm{End}\,E_\varepsilon)
\fe
vanishes. This is the Atiyah/holomorphy constraint \cite{Atiyah:1957xx,Anderson:2011ty}: complex-structure deformations survive only in the kernel of $\alpha_{E_\varepsilon}$. At the standard embedding point $\alpha_{E_0}=0$ for all $[Z]$ (equivalently, $\iota_Z R$ is $\bar\partial$-exact, as encoded by (\ref{eq:Lambda_tilde})), but for a generic nearby bundle deformation $\alpha_{E_\varepsilon}$ becomes nontrivial. Expanding ${\cal F}_\varepsilon=R+\varepsilon\,\delta{\cal F}+O(\varepsilon^2)$, the first obstruction to moving simultaneously in complex structure is therefore the \emph{linearized} constraint
\ie\label{eq:Atiyah_lin}
[\iota_Z\,\delta{\cal F}] \;=\; 0 \qquad \text{in } H^2(X,\mathrm{End}\,T_X),
\fe
which depends on the chosen direction away from $E=T_X$. Equivalently, the simultaneous first-order deformations of $(X,E_\varepsilon)$ are packaged by the extension bundle $Q_\varepsilon$ defined by \cite{Atiyah:1957xx,Anderson:2011ty}
\ie\label{eq:Atiyah_extension}
0 \;\to\; \mathrm{End}\,E_\varepsilon \;\to\; Q_\varepsilon \;\to\; T_X \;\to\; 0,
\fe
and the long exact sequence implies that the image of $H^1(X,Q_\varepsilon)$ in $H^1(X,T_X)$ is precisely $\ker(\alpha_{E_\varepsilon})$.

\subsection{Interpreting the deformation classes of section 3.1}

Sections 4.1--4.2 describe which infinitesimal deformations exist as right-chiral operators in the large-volume NLSM and how their \emph{linear} compatibility conditions are encoded by the Atiyah map. For exact marginality, one must additionally require that turning on such deformations does not generate a beta function. At leading order this is controlled by the NS--NS--NS correlators with a PCO insertion discussed in sections 3.2--3.3; at large volume these correlators reduce to the deformation-theory brackets summarized at the end of section 3.4. We now reinterpret the operator classes of section 3.1 in this geometric language, emphasizing which NS--NS--NS correlators detect their first obstructions.

\subsubsection*{Gauge-neutral directions: $(Z,Y,\widetilde\Lambda)$}

At the standard embedding point, gauge-neutral supersymmetric marginal operators split as
\ie
\relax [Z]\in H^1(X,T_X),\qquad [Y]\in H^1(X,T_X^\vee),\qquad [\widetilde\Lambda]\in H^1(X,\mathrm{End}\,T_X),
\fe
with $Z,Y$ corresponding to the $q=0$ descendant moduli (\ref{defmoduli}) and $\widetilde\Lambda$ corresponding to the current-type multiplets (\ref{defcurrent}). The linear Melnikov--Sharpe conditions (\ref{eq:geom_cond}) ensure that these deformations exist at first order. The leading CPT constraints come from the quadratic brackets:
\begin{enumerate}
\item {\bf Complex structure ($Z$): integrability.}\\
For a complex-structure deformation to extend beyond first order one must impose integrability, schematically
\ie
\bar\partial Z \;+\; \frac12\,[Z,Z] \;=\; 0,
\fe
where $[Z,Z]\in\Omega^{0,2}(X,T_X)$ is the Kodaira--Spencer bracket of section 3.4. The corresponding obstruction class lies in $H^2(X,T_X)$, and its projection onto marginal directions is captured by an NS--NS--NS correlator of the form (\ref{NSNSNS3pt}) built from the $Z$ vertices. For a smooth Calabi--Yau threefold, the standard Tian--Todorov argument implies that this obstruction vanishes, in accord with the existence of the $(2,2)$ conformal manifold.

\item {\bf Bundle deformations away from the standard embedding ($\widetilde\Lambda$): holomorphy and self-obstructions.}\\
Turning on $\widetilde\Lambda$ deforms the holomorphic structure on $E\simeq T_X$; in components this is the deformation of the $(0,1)$ connection $A^{0,1}\to A^{0,1}+\widetilde\Lambda$. The condition that the deformed bundle remain holomorphic is the Maurer--Cartan equation
\ie\label{eq:MC_bundle}
\bar\partial\,\widetilde\Lambda \;+\; \frac12\,[\widetilde\Lambda,\widetilde\Lambda] \;=\; 0,
\fe
where $[\widetilde\Lambda,\widetilde\Lambda]\in\Omega^{0,2}(X,\mathrm{End}\,T_X)$ is the bundle bracket of section 3.4. The corresponding obstruction class lies in $H^2(X,\mathrm{End}\,T_X)$; when nontrivial it obstructs continuing the $\widetilde\Lambda$ deformation beyond first order. In the worldsheet language this is precisely what is measured by the NS--NS--NS correlators involving three current-type multiplets in section 3.3: those correlators compute the OPE coefficients that enter the quadratic beta function for the $\widetilde\Lambda$ couplings.

\item {\bf Mixed $(Z,\widetilde\Lambda)$ constraints: the Atiyah map.}\\
In a background where the bundle is deformed to $E_\varepsilon$ (equivalently $\langle\widetilde\Lambda\rangle\neq 0$), the would-be complex-structure modulus $Z$ must also remain compatible with bundle holomorphy, i.e.\ satisfy $\alpha_{E_\varepsilon}([Z])=0$ as in (\ref{eq:Atiyah_map}). Linearizing around the standard embedding gives the bilinear constraint (\ref{eq:Atiyah_lin}), which is a class in $H^2(X,\mathrm{End}\,T_X)$. In conformal perturbation theory this constraint is encoded in the mixed OPE coefficient, equivalently by the sphere NS--NS--NS correlator of integrated marginal operators (with one PCO insertion, cf.\ sections 3.2--3.3)
\ie\label{eq:Atiyah_CPT_3pt}
C^{a}{}_{Ib}\;\propto\;\left\langle \Phi^{a}(\infty)\,\Phi^{I}(1)\,\Phi^{b}(0)\right\rangle,
\qquad
\Phi^{i}\equiv \wt G^-_{-{1\over 2}}{\cal W}_i,
\fe
which is the natural mixed OPE coefficient that would govern the CFT avatar of the Atiyah/holomorphy constraint in a bundle-deformed background.
At the strict standard embedding $(2,2)$ point, however, the selection rules of section 3.3 imply $C^{a}{}_{Ib}=0$ whenever $\Phi^{I}$ is a $q=0$ descendant modulus (\ref{defmoduli}) and $\Phi^{a},\Phi^{b}$ are current-type: the gauge-fermion sector contributes a nonzero $\langle J\,J\rangle$, but the \emph{internal} holomorphic sector sees only a single nontrivial insertion because the current-type operators lie in the left vacuum module. Thus the internal holomorphic factor reduces to a one-point function (equivalently a vacuum matrix element), which vanishes on the sphere for a marginal deformation (or, more explicitly for $q=0$ descendants, by the holomorphic ${\cal N}=2$ Ward identity of section 3.2).
Concretely, using the unintegrated representatives of section 3.2, one would extract $C^{a}{}_{Ib}$ from a correlator of the form
\[
C^{a}{}_{Ib}\;\propto\;
\left\langle
V_{\delta A_a}^{(0)}(\infty)\,
V_{\mu_I}^{(-1)}(1)\,
\overline V_{\delta A_b}^{(-1)}(0)
\right\rangle,
\]
with $V_{\delta A}^{(-1)}\sim e^{-{\wt\phi}}J^{A'B'}\delta A_{\bar i}^{A'B'}(X)\,\wt\psi^{\bar i}$ and $V_{\mu}^{(-1)}\sim e^{-{\wt\phi}}\,G^+_{-{1\over 2}}\big(\mu_{\bar i}{}^{j}(X)\,\psi_{j}\,\wt\psi^{\bar i}\big)$ as in section 3.4, and $V^{(0)}=X\cdot V^{(-1)}$ obtained by picture-changing; the correlator vanishes at the strict $(2,2)$ point by the same holomorphic one-point selection rule.

The linearized Atiyah lifting is therefore first visible in perturbation theory through the integrated four-point functions that correct this mixed OPE data once the current-type background is switched on, see section \ref{sec:4pt_selection} (and (\ref{eq:Atiyah_4pt_LV}) for a representative large-volume form, reducing to $[\iota_Z\,\delta{\cal F}]$ in (\ref{eq:Atiyah_lin})).
\end{enumerate}

The $Y$ directions (K\"ahler and $B$-field moduli) also enter the full $(0,2)$ deformation problem, but for the tree-level $H=0$ backgrounds we consider here they decouple from the holomorphy/Atiyah constraint at linear order; incorporating the $\alpha'$-corrected anomaly/Bianchi constraints would re-couple $Y$ to $(Z,\widetilde\Lambda)$ and is beyond our present scope.

\subsubsection*{Charged and model-dependent directions}

The remaining deformation classes of section 3.1 are not part of the gauge-neutral $(Z,Y,\widetilde\Lambda)$ system, but their leading obstructions are governed by the same NS--NS--NS/CPT logic.
\begin{enumerate}
\item {\bf Matter/Higgs directions ($\Lambda=\lambda$).}\\
Giving a vev to a charged matter multiplet deforms the bundle by turning on off-diagonal components of the connection, often described by extension data. The resulting compatibility constraints on complex structure can be understood as Atiyah-type conditions for the corresponding deformed bundle, and the relevant CPT data are again NS--NS--NS correlators with one insertion of a complex-structure operator and two insertions in the Higgs/background direction. A key effect emphasized in \cite{Anderson:2010mh,Anderson:2011ty} is that such extension classes can exist only on special loci in complex-structure moduli space (``jumping'' cohomology).
\item {\bf Spectral-flow BPS multiplets ($q=\pm 2$, $\Lambda=\mathbf{1}$).}\\
At large volume these are represented by operators valued in $H^1(X,\wedge^2 T_X)$ and $H^1(X,\wedge^2 T_X^\vee)$ (section 3.4). Their quadratic obstructions are governed by the corresponding brackets/cup products into $H^2(X,\wedge^2 T_X)$ or $H^2(X,\wedge^2 T_X^\vee)$, and the associated OPE coefficients are extracted from the NS--NS--NS correlators listed in section 3.3.
\item {\bf Additional $q=0$ primaries (long multiplets).}\\
These require extra holomorphic weight-one operators beyond the ${\cal N}=2$ algebra and are absent at a generic smooth Calabi--Yau large-volume point, but can arise at special free-field/orbifold points. When present, their exact marginality is tested in the same way: by the vanishing of the corresponding CPT OPE coefficients computed by NS--NS--NS correlators with a PCO insertion.
\end{enumerate}

\section{Conformal perturbation theory and obstructions}\label{sec:cpt_obstructions}

We now phrase the obstruction problem directly in conformal perturbation theory. Let $\Phi_i$ be a basis of marginal operators of the undeformed theory (in the present setting, $\Phi_i$ are the integrated versions of the supersymmetric deformations (\ref{rFterm})). Consider the perturbation
\ie\label{cptpert}
S \;\to\; S + \sum_i \lambda^i \int d^2z~\Phi_i(z,\bz).
\fe
In conformal perturbation theory, the beta functions for the couplings $\lambda^i$ are determined by logarithmic divergences in integrated correlators, equivalently by the singular terms in operator product expansions. In a scheme where the $\Phi_i$ are marginal at $\lambda=0$ and orthonormalized with respect to the two-point function, the leading term takes the universal form
\ie\label{beta2}
\beta^i \;=\; \pi\, C^i{}_{jk}\,\lambda^j\lambda^k + O(\lambda^3),
\fe
where the coefficients $C^i{}_{jk}$ appear in the OPE
\ie\label{OPE}
\Phi_j(z,\bz)\,\Phi_k(0)\;\sim\; {1\over |z|^2}\,C^i{}_{jk}\,\Phi_i(0) + \cdots.
\fe
Exact marginality requires that all beta functions vanish order-by-order; a first obstruction is therefore a nonzero projection of the OPE (\ref{OPE}) back onto the space of marginal operators. A second kind of obstruction arises if the OPE produces relevant operators: such couplings are then generated along the RG flow and must be tuned away, which is impossible if they violate the symmetries one insists on preserving.

\subsection*{Summary of NS obstructions and physical meaning}

An ``NS obstruction'' is a failure of exact marginality in conformal perturbation theory: a nonzero OPE projection back onto marginal operators generates a beta function, and if the leading (quadratic) coefficients vanish by symmetry, the first obstruction can move to higher order and is governed by integrated higher-point functions.
\subsubsection*{Pure ``bundle'' deformations}
Take couplings $\varepsilon^a$ for current-type marginal operators $\Phi_a$ (section 3.1), \emph{i.e.}\ current--current insertions after picture-changing. Their leading NS obstruction is a nonzero current-type OPE coefficient $C^{a}{}_{bc}$, which induces $\beta^a\propto C^{a}{}_{bc}\,\varepsilon^b\varepsilon^c$. A protected unobstructed block is obtained by restricting to commuting currents on each side (Chaudhuri--Schwartz), in which case the current--current perturbation is integrably marginal. In a large-volume sigma-model limit these $\Phi_a$ reduce to the usual gauge-neutral ``bundle'' directions (section 3.4).

\subsubsection*{Mixed complex-structure--bundle deformations}
Let $\Phi_I$ denote the $q=0$ descendant moduli (the $(2,2)$ conformal-manifold directions, interpreted as complex-structure-type at large volume) and $\Phi_b$ a current-type operator. At the strict $(2,2)$ point, the gauge-fermion factor $\langle J\,J\rangle$ is nonzero, but the internal holomorphic piece reduces to a one-point function, so the would-be mixed NS--NS--NS coefficient with two current-type insertions vanishes by the holomorphic one-point selection rule (section 3.3); the first mixed obstructions are therefore controlled by integrated four-point functions once the current-type background is turned on (section \ref{sec:4pt_selection}). In the large-volume limit this reduces to the usual complex-structure/bundle compatibility (Atiyah-type) condition of section 3.4.

\subsubsection*{Yukawas and the spacetime potential}
NS--R--R three-point functions compute the cubic superpotential couplings $Y_{ijk}$ (section 2 and (\ref{yukawa3pt})), hence the leading F-term potential generated in spacetime. They are related by spectral flow to BPS NS data, but exact marginality is decided by the NS--NS--NS (and higher-point) correlators of the integrated perturbing operators that enter the worldsheet beta functions.

\subsection{Selection rules from ${\cal N}=2$ symmetry at the standard embedding point}

At the standard embedding point the internal theory has a holomorphic ${\cal N}=2$ SCA, so holomorphic $U(1)_R$ charge and Ward identities impose strong constraints on the OPEs of the candidate deformations classified in section 3.1. In particular:
\begin{itemize}
\item For the $(2,2)$-preserving moduli (the $q=0$ descendants), the perturbing operators sit in short multiplets built on chiral primaries. Their OPEs are constrained by the chiral ring, and the corresponding directions are expected to be exactly marginal, giving the usual conformal manifold of the $(2,2)$ theory.
	\item For deformations that leave the standard embedding locus (those involving $\lambda^{A'}$, the gauge currents $J^{A'B'}$, or carrying $q=\pm 2$), the holomorphic ${\cal N}=2$ symmetry is broken. In conformal perturbation theory this means that the cancellations implied by ${\cal N}=2$ Ward identities need not persist, and one generically expects nonzero coefficients $C^i{}_{jk}$ in (\ref{beta2}) unless forbidden by remaining global symmetries or charge selection rules.
		\end{itemize}
			In practice, one can proceed by starting with a candidate linear combination of deformations from section 3.1, using charge conservation to determine which OPE channels are allowed, and then computing (or constraining) the corresponding three- and four-point functions that control the first nontrivial terms in the beta functions. The vanishing of the cubic superpotential couplings computed in section 2 is a necessary condition for F-flatness, but exact marginality requires the absence of logarithmic divergences to all orders in conformal perturbation theory.

			\subsubsection*{Ward-identity constraints on NS--NS--NS CPT coefficients}
			
			At the standard embedding point the holomorphic ${\cal N}=2$ SCA is conserved, so its Ward identities constrain not only the NS--R--R Yukawas of section 3.2 but also the NS--NS--NS correlators (\ref{NSNSNS3pt}) that determine the quadratic CPT coefficients. Since picture-changing inserts only the \emph{right-moving} supercurrent (\ref{V0fromPCO}), the holomorphic Ward-identity manipulations are unchanged.
			
			The basic ingredient is the same contour argument used in section 3.2. For a correlator with one holomorphic super-descendant of a chiral primary, of the schematic form (\ref{Gdesc3pt}),
			\[
			\left\langle \big(G_{-{1\over 2}}^\pm \phi\big)(\infty)\,\psi(1)\,\chi(0)\right\rangle,
			\]
			one represents the insertion at infinity as a contour integral $\oint {dz\over 2\pi i}\,z\,G^\pm(z)$ and shrinks the contour. If $\psi$ and $\chi$ are holomorphic primaries, the term where the positive mode $G^\pm_{1\over 2}$ hits the insertion at the origin vanishes, and one is left with a correlator where $G^\pm_{-{1\over 2}}$ acts on the remaining insertion. In particular, if that remaining insertion is a chiral (respectively anti-chiral) primary annihilated by $G^+_{-{1\over 2}}$ (respectively $G^-_{-{1\over 2}}$), the correlator vanishes.
			
			Applied to the deformation multiplets of section 3.1, this gives an extra model-independent vanishing criterion at the $(2,2)$ point: any NS--NS--NS coefficient in which the $q=0$ insertion is a standard-embedding modulus of descendant form (\ref{defmoduli}) and the other two holomorphic insertions are BPS primaries (the $q=\pm 1$ matter multiplets (\ref{defmatter}) or the $q=\pm 2$ spectral-flow multiplets (\ref{defq2})) vanishes. Equivalently, within the list of ``possibly non-vanishing'' NS--NS--NS correlators in section 3.3, the families with holomorphic charge patterns $(+1,-1,0)$ and $(0,+2,-2)$ receive contributions only if the $q=0$ insertion is a $q=0$ \emph{primary} (\ref{defq0primary}) rather than a descendant modulus.
			
			By contrast, correlators in the $(0,0,0)$ family are not constrained by this Ward identity, since the remaining $q=0$ insertions need not be chiral/anti-chiral primaries. Moreover, once one perturbs away from the standard embedding locus, the holomorphic ${\cal N}=2$ current is no longer conserved, so the Ward-identity vanishings above need not persist: the corresponding obstructions can reappear through higher-order CPT data (integrated four-point functions), cf.\ section \ref{sec:4pt_selection}.
			
				\subsubsection*{Integrably marginal current--current bundle deformations}
				
				While a generic current-type (``bundle'') deformation is expected to be obstructed by nonzero CPT coefficients, there is an important protected subclass that can be identified purely at the level of current-algebra data in the worldsheet CFT.
				Suppose the theory contains a set of holomorphic currents $J^a(z)$ and anti-holomorphic currents $\bar J^{\bar b}(\bz)$ generating (possibly distinct) affine current algebras. Consider a deformation by a current--current operator
			\[
			\delta S \;=\; \sum_{a,\bar b} g_{a\bar b}\int d^2z~J^a(z)\,\bar J^{\bar b}(\bz),
			\]
			which is manifestly of weight $(1,1)$. The key point is that the leading CPT obstruction is controlled directly by the single-pole terms in the current OPEs. Writing
			\[
			J^a(z)\,J^c(0)\sim {k^{ac}\over z^2}+{f^{ac}{}_{e}\,J^e(0)\over z}+\cdots,
			\qquad
			\bar J^{\bar b}(\bz)\,\bar J^{\bar d}(0)\sim {\bar k^{\bar b\bar d}\over \bz^2}+{\bar f^{\bar b\bar d}{}_{\bar f}\,\bar J^{\bar f}(0)\over \bz}+\cdots,
			\]
			one finds for the marginal operators ${\cal O}_{a\bar b}\equiv J^a\bar J^{\bar b}$ the schematic closure
			\[
			{\cal O}_{a\bar b}(z,\bz)\,{\cal O}_{c\bar d}(0)\;\sim\;{1\over |z|^2}\,f^{ac}{}_{e}\,\bar f^{\bar b\bar d}{}_{\bar f}\,{\cal O}_{e\bar f}(0)+\cdots,
			\]
			so the quadratic beta function for $g_{a\bar b}$ is proportional to the product of structure constants. A sufficient criterion for such a perturbation to be \emph{integrably marginal} (hence exactly marginal to all orders in conformal perturbation theory) was given by Chaudhuri and Schwartz \cite{Chaudhuri:1988qb}. In particular, if the currents used in the deformation generate an \emph{abelian} subalgebra on each side (so the corresponding OPEs have no single-pole terms), then the deformation is integrably marginal: no logarithmic divergences are generated at any order and the beta functions vanish identically.
			
			This also gives a purely CFT lower bound on unobstructed current-type deformations: if one can identify commuting sets of currents $\{J^{a}\}_{a=1}^{r_L}$ and $\{\bar J^{\bar b}\}_{\bar b=1}^{r_R}$, then the $r_L r_R$ couplings $g_{a\bar b}$ in that block parameterize an exactly marginal locus. For a compact semisimple current algebra one may always take $r_{L,R}$ to be the rank, i.e.\ the dimension of a Cartan subalgebra.
			
			This criterion applies directly to a subset of the current-type directions in our classification. Recall that a current-type multiplet has left factor a gauge current $J^{A'B'}$ and a right-moving chiral operator ${\cal O}^{\rm NS}$ in the left vacuum module. After picture-changing, the integrated insertion involves
			\[
			V^{(0)} \;\sim\; J^{A'B'}\,\wt G^-_{-{1\over 2}}{\cal O}^{\rm NS},
			\]
			and since ${\cal O}^{\rm NS}$ lies in the left vacuum module, $\wt G^-_{-{1\over 2}}{\cal O}^{\rm NS}$ is an anti-holomorphic weight-one operator, i.e.\ a conserved current $\bar j(\bz)$. Thus current-type directions are of current--current form, and the Chaudhuri--Schwartz criterion guarantees exact marginality whenever one restricts to an abelian set of left and right currents. Concretely, at free-field points one may take ${\cal O}^{\rm NS}=\wt\psi^{\bar i}$ so that $\wt G^-_{-{1\over 2}}\wt\psi^{\bar i}\sim \bar\partial X^{\bar i}$, and choose the left-moving current from a commuting Cartan subalgebra of the gauge current algebra; the resulting deformation $\int d^2z~J^{\rm Cartan}(z)\,\bar\partial X^{\bar i}(\bz)$ is then exactly marginal by the current--current criterion.
		
		\subsection{Four-point functions and subleading selection rules}\label{sec:4pt_selection}
		
		At leading order, the coefficients $C^i{}_{jk}$ in (\ref{beta2}) are extracted from NS--NS--NS correlators with one PCO insertion, cf.\ (\ref{NSNSNS3pt}). At the next order in conformal perturbation theory, potential logarithmic divergences are governed by integrated four-point functions of the marginal operators $\Phi_i$. Since the perturbation (\ref{cptpert}) involves only the integrated NS operators $\Phi_i$, the relevant worldsheet data are sphere NS--NS--NS--NS correlators with two PCO insertions; correlators with Ramond insertions compute spacetime couplings but do not enter the beta functions for (\ref{cptpert}).
	
	On the sphere the total right-moving picture number must be $-2$, so a four-point function requires two PCO insertions. Choosing the picture assignment $(0,0,-1,-1)$, a representative correlator is
	\ie\label{NSNSNSNS4pt}
	\mc{I}_{ij k\bar\ell}\;\propto\;\int d^2z~\left\langle V_i^{(0)}(\infty)\,V_j^{(0)}(1)\,V_k^{(-1)}(z)\,\overline V_\ell^{(-1)}(0)\right\rangle,
	\fe
	together with permutations and the conjugate correlators with $V\leftrightarrow \overline V$, and variants where one or both 0-picture vertices are $\overline V^{(0)}$. Here $V^{(0)}$ and $\overline V^{(0)}$ are obtained by picture-changing as in (\ref{V0fromPCO}), so (\ref{NSNSNSNS4pt}) contains two explicit insertions of the right-moving supercurrent.
	
	The same three selection rules as in section 3.2 extend immediately to four points:
	\begin{enumerate}
	\item {\bf Right-moving $U(1)_R$ (picture/PCO).}\\
	Since the 0-picture vertices have $\wt q=0$, right-charge conservation implies that the two $(-1)$-picture vertices must have opposite right charges. In particular, in (\ref{NSNSNSNS4pt}) one must choose one $(-1)$-picture insertion from the $\wt q=+1$ sector and one from the $\wt q=-1$ sector (cf.\ (\ref{qtildeSelection})).
	
	\item {\bf Gauge-fermion Wick contractions.}\\
	The free $\lambda^{A'}$ system again factorizes, and Wick's theorem restricts the possible gauge structures. Up to trivial identity insertions, the only potentially nonzero correlators built from $\Lambda\in\{\mathbf{1},\lambda^{A'},J^{A'B'}\}$ are of the schematic types
	\ie\label{gaugecontractions4pt}
	&\langle \lambda\,\lambda\rangle\neq 0,\qquad
	\langle \lambda\,\lambda\,\lambda\,\lambda\rangle\neq 0,\qquad
	\langle J\,\lambda\,\lambda\rangle\neq 0,\qquad
	\langle J\,J\rangle\neq 0,
	\\
	&\langle J\,J\,\lambda\,\lambda\rangle\neq 0,\qquad
	\langle J\,J\,J\rangle\neq 0,\qquad
	\langle J\,J\,J\,J\rangle\neq 0.
	\fe
	In particular, a single $J$ insertion with no external $\lambda$'s vanishes, while any correlator with an odd number of single-$\lambda$ insertions vanishes.
	
	\item {\bf Holomorphic $U(1)_R$ charge.}\\
	Conservation of the holomorphic $U(1)_R$ current implies
	\ie\label{qsum0_4pt}
	q_1+q_2+q_3+q_4=0,
	\fe
	with $q\in\{0,\pm 1,\pm 2\}$ for the multiplets of section 3.1. Since the $q=\pm 1$ multiplets are precisely the single-$\lambda$ matter directions, the gauge selection rule forces the number of $\pm 1$ insertions to be even. Up to permutation, the only allowed holomorphic charge patterns are therefore
	\ie
	&(q,q,q,q)=(+1,+1,-1,-1),
	\\
	&(q,q,q,q)=(+1,-1,0,0),\quad (+1,-1,+2,-2),\quad (+1,+1,-2,0),\quad (-1,-1,+2,0),
	\\
	&(q,q,q,q)=(0,0,0,0),\quad (0,0,+2,-2),\quad (+2,+2,-2,-2).
	\fe
	\end{enumerate}
	
	Combining these rules, the NS--NS--NS--NS correlators among the deformation multiplets that are not forced to vanish at a generic standard-embedding $(2,2)$ point are:
	\begin{enumerate}
	\item {\bf Purely internal ($\Lambda=\mathbf{1}$ for all four).}\\
	The gauge factor is trivial. Holomorphic charge conservation allows
	\[
	(0,0,0,0),\qquad (0,0,+2,-2),\qquad (+2,+2,-2,-2),
	\]
	with the $q=0$ sector including both descendants (\ref{defmoduli}) and possible primaries (\ref{defq0primary}).
	
	\item {\bf Two matter and two $F=0$ insertions ($\langle\lambda\,\lambda\rangle\neq 0$).}\\
	With two single-$\lambda$ multiplets, (\ref{qsum0_4pt}) allows the holomorphic patterns $(+1,-1,0,0)$,\allowbreak $(+1,-1,+2,-2)$,\allowbreak $(+1,+1,-2,0)$,\allowbreak and $(-1,-1,+2,0)$.
	
	\item {\bf Four matter insertions ($\langle\lambda\,\lambda\,\lambda\,\lambda\rangle\neq 0$).}\\
	Holomorphic charge conservation forces two $q=+1$ and two $q=-1$ matter multiplets, i.e.\ $(+1,+1,-1,-1)$.
	
	\item {\bf One current, two matter, one $F=0$ insertion ($\langle J\,\lambda\,\lambda\rangle\neq 0$).}\\
	The current carries $q=0$, so (\ref{qsum0_4pt}) permits
	\[
	(+1,-1,0,0),\qquad (+1,+1,-2,0),\qquad (-1,-1,+2,0).
	\]
	
	\item {\bf Two currents and two matter ($\langle J\,J\,\lambda\,\lambda\rangle\neq 0$).}\\
	Holomorphic charge conservation forces the two matter multiplets to have opposite charge, giving $(+1,-1,0,0)$.
	
	\item {\bf Two currents and two $F=0$ insertions ($\langle J\,J\rangle\neq 0$).}\\
	The $F=0$ charges must sum to zero, allowing $(0,0,0,0)$ and $(0,0,+2,-2)$.
	
	\item {\bf Three currents and one $q=0$ $F=0$ insertion ($\langle J\,J\,J\rangle\neq 0$).}\\
	All holomorphic charges are zero, so the remaining insertion must lie in the $q=0$ sector.
	
	\item {\bf Four currents ($\langle J\,J\,J\,J\rangle\neq 0$).}\\
	Holomorphic charge conservation is automatic (all $q=0$); the only model-independent constraint is the right-moving picture/charge condition above.
	\end{enumerate}
	In every case one must also satisfy the right-moving selection rule: in a fixed $(0,0,-1,-1)$ picture assignment, one of the two $(-1)$-picture insertions is chosen from the $\wt q=-1$ sector and the other from the $\wt q=+1$ sector, as in (\ref{qtildeSelection}).
	
		\subsubsection{Large-volume interpretation I: mixed Atiyah lifting from four-point data}
		
		At large volume, the vertices of section 3.4 have free-field representatives built from target-space tensors contracted with the right-moving fermions $\wt\psi^{\bar i}$. In an NS--NS--NS--NS correlator, the two PCO insertions produce two 0-picture vertices via (\ref{V0fromPCO}); at leading order this implements the replacement $\wt\psi^{\bar i}\to \bar\partial X^{\bar i}$ as in (\ref{GpsiX}). As a result, the 4-point data that enter CPT can often be read geometrically as higher-order terms in the deformation-theory maps for complex structure and bundles (Kuranishi/Atiyah-type constraints), even when the leading 3-point coefficient vanishes at the strict $(2,2)$ point.
		
		A particularly important example is the lifting of complex-structure directions by the Atiyah constraint. At the strict standard embedding $(2,2)$ point the corresponding mixed two-current NS--NS--NS coefficient vanishes (cf.\ (\ref{eq:Atiyah_CPT_3pt})), in accord with $\alpha_{E_0}=0$. The first nontrivial lifting at small current-type background is therefore governed by integrated four-point functions. For instance, the leading mixed coefficient induced linearly in a current-type background direction $\delta A_\varepsilon\in H^{1}(X,\mathrm{End}\,T_X)$ is schematically extracted from a correlator in the ``three currents + one $q=0$'' family,
\ie\label{eq:Atiyah_4pt_LV}
\delta_\varepsilon C^{a}{}_{Ib}
\;\propto\;
\int d^2z~\left\langle
	V_{\delta A_a}^{(0)}(\infty)\,
	V_{\delta A_\varepsilon}^{(0)}(1)\,
	V_{\mu_I}^{(-1)}(z)\,
	\overline V_{\delta A_b}^{(-1)}(0)
	\right\rangle,
	\fe
	where $V_{\delta A}^{(-1)}\sim e^{-{\wt\phi}}\,J^{A'B'}\,\delta A_{\bar i}^{A'B'}(X)\,\wt\psi^{\bar i}$ is the large-volume representative of a current-type vertex (section 3.4), $V^{(0)}=X\cdot V^{(-1)}$ is obtained by picture-changing, and the complex-structure modulus insertion is of descendant form
	\[
	V_{\mu}^{(-1)}\;\sim\; e^{-{\wt\phi}}\,G^+_{-{1\over 2}}\big(\mu_{\bar i}{}^{k}(X)\,\psi_{k}\,\wt\psi^{\bar i}\big)
	\;\sim\;
	e^{-{\wt\phi}}\,\mu_{\bar i}{}^{k}(X)\,\partial X^{k}\,\wt\psi^{\bar i}+\cdots.
	\]
	To see explicitly how (\ref{eq:Atiyah_4pt_LV}) reproduces the usual Atiyah class, it is convenient to work in a local holomorphic gauge $A_k=0$ for the background bundle connection, so that the $(1,1)$ curvature is ${\cal F}_{\bar j k}=\partial_k A_{\bar j}$ and the linearized curvature variation induced by $\delta A_\varepsilon$ is
	\ie\label{eq:deltaF_hol_gauge}
	\delta{\cal F}(\delta A_\varepsilon)_{\bar j k}\;=\;\partial_k\,\delta A_{\varepsilon\bar j}.
	\fe
	Using the free-field OPE $X^{k}(z)\,X^{\ell}(w)\sim -g^{k\ell}\log|z-w|^2$ gives
	\[
	\partial X^{k}(z)\,\delta A_{\varepsilon\bar j}(X(1))
	\;\sim\;
	-{1\over z-1}\,\partial_k \delta A_{\varepsilon\bar j}(X(1))+\cdots,
	\]
	so the short-distance OPE of the complex-structure vertex with the background current-type insertion produces a singular term proportional to
	\[
	V_{\mu_I}^{(-1)}(z)\,V_{\delta A_\varepsilon}^{(0)}(1)
	\;\sim\;
	-{1\over z-1}\;
	e^{-{\wt\phi}}\,J^{A'B'}\,
	\big(\iota_{\mu_I}\delta{\cal F}(\delta A_\varepsilon)\big)_{\bar i\bar j}^{A'B'}(X(1))\,
	\wt\psi^{\bar i}\,\bar\partial X^{\bar j}
	+\cdots,
	\]
	where $(\iota_{\mu}\delta{\cal F})_{\bar i\bar j}\equiv \mu_{[\bar i}{}^{k}\,\delta{\cal F}_{\bar j]k}$ as in (\ref{eq:Atiyah_lin}). Inserting this OPE into (\ref{eq:Atiyah_4pt_LV}) shows that the first nonzero correction to the mixed OPE data is governed precisely by the cohomology class
	\[
	[\iota_{Z_I}\,\delta{\cal F}(\delta A_\varepsilon)]\in H^{2}(X,\mathrm{End}\,T_X)
	\]
		onto the current-type marginal directions. Requiring that this correction vanish for the complex-structure deformation $Z_I$ is therefore equivalent (at leading order in the bundle-background parameter $\varepsilon$) to the linearized Atiyah/holomorphy constraint (\ref{eq:Atiyah_lin}).
		
		\subsubsection{Large-volume interpretation II: pure bundle obstructions and $C^{a}{}_{bc}$}
		
		For a \emph{pure} current-type (``bundle'') deformation with couplings $\varepsilon^a$, the leading CPT obstruction is the quadratic beta function
		\[
		\beta^a\;=\;\pi\,C^{a}{}_{bc}\,\varepsilon^b\varepsilon^c+O(\varepsilon^3),
		\]
			where $C^{a}{}_{bc}$ is extracted from the NS--NS--NS correlators with three current-type multiplets in section 3.3. In the large-volume sigma-model realization, these multiplets are represented by $(0,1)$ bundle fluctuations $\delta A_a\in\Omega^{0,1}(X,\mathrm{End}\,T_X)$, and the corresponding integrated marginal operators take the schematic form
		\[
		\Phi_{\delta A_a}\;\sim\;\delta A_{a\bar i}^{A'B'}(X)\,\bar\partial X^{\bar i}\,J^{A'B'}+\cdots,
		\]
			as in (\ref{rFterm}) and (\ref{GpsiX}). The same free-field reasoning as above then shows that the short-distance OPE coefficients are controlled by the bundle bracket (3.42): the commutator in $\mathrm{End}\,T_X$ comes from the single-pole term in the $J\,J$ OPE, and antisymmetrization in the $(0,1)$ indices is inherited from the two $\wt\psi^{\bar i}$ insertions. Equivalently, by Serre duality on a Calabi--Yau threefold, the sphere three-point function with three current-type insertions pairs $[\delta A_a]\in H^{1}(X,\mathrm{End}\,T_X)$ with the obstruction class $[[\delta A_b,\delta A_c]]\in H^{2}(X,\mathrm{End}\,T_X)$, so $C^{a}{}_{bc}$ is generically nonzero unless the commutator bracket is cohomologically trivial. In particular, if the bundle fluctuations take values in a commuting (abelian) subalgebra, then $[\delta A_b,\delta A_c]=0$ and the corresponding block has vanishing $C^{a}{}_{bc}$, in accord with the Chaudhuri--Schwartz integrably marginal criterion discussed above.
		
		More generally, the allowed correlators with three or four current-type insertions encode higher-order terms in the Kuranishi map for deformations of the holomorphic structure on $E\simeq T_X$. In geometric deformation theory these higher obstructions are expressed in terms of iterated brackets and Green's operators (Massey/Yoneda-type products) valued in $H^{2}(X,\mathrm{End}\,T_X)$; in CPT they arise as the $O(\lambda^3)$ contributions to the beta functions for the bundle couplings and are governed precisely by the corresponding NS--NS--NS--NS correlators with two PCO insertions.
	
		Similarly, the non-vanishing selection-rule families involving matter/Higgs insertions ($\Lambda=\lambda^{A'}$) admit a large-volume interpretation in terms of bundle-valued cohomology and its variation. For instance, a matter fluctuation is represented by a class $[\alpha]\in H^{1}(X,V_R)$ for an associated bundle $V_R$. Away from the strict $(2,2)$ locus, complex-structure and bundle deformations vary the Dolbeault operator $\bar\partial_{V_R}$; at the linearized level this is encoded by the associated-bundle Atiyah map
	\ie
	\alpha_{V_R}([Z])([\alpha]) \;\equiv\; [\iota_Z{\cal F}_{V_R}\cdot \alpha]\;\in\; H^{2}(X,V_R),
	\fe
	where ${\cal F}_{V_R}$ is the curvature of the induced connection on $V_R$. When nontrivial, such classes obstruct continuing the corresponding matter direction and are the geometric origin of ``jumping'' phenomena. Higher-order compatibility conditions are expressed by iterated Yoneda/Massey products among the extension classes, and their CFT avatars appear as the allowed 4-point correlators with two or four single-$\lambda$ insertions listed above.
	
	
	\bibliographystyle{JHEP}
	\bibliography{stringrefs}


\end{document}
